\documentclass[10pt,reqno]{amsart}

\usepackage[T1]{fontenc}
\usepackage{lmodern}
\usepackage{microtype}
\usepackage{amsmath,amssymb,amsthm,mathtools}
\usepackage{booktabs,longtable,array,tabularx}
\usepackage{enumitem}
\usepackage{xcolor}
\usepackage{listings}
\usepackage{seqsplit}
\usepackage{geometry}
\usepackage{hyperref}
\usepackage[nameinlink,noabbrev]{cleveref}

\geometry{letterpaper,margin=1in}
\hypersetup{
  colorlinks=true,
  linkcolor=blue!45!black,
  citecolor=green!35!black,
  urlcolor=blue!55!black,
  pdftitle={An explicit degree-494 candidate class for the Cappell--Shaneson sphere X(41,189,73)},
  pdfauthor={},
  pdfsubject={A conditional skein-lasagna obstruction with exact finite computations}
}

\setlist[itemize]{leftmargin=2.1em,itemsep=2pt,topsep=4pt}
\setlist[enumerate]{leftmargin=2.3em,itemsep=2pt,topsep=4pt}
\lstset{
  basicstyle=\ttfamily\small,
  columns=fullflexible,
  breaklines=true,
  frame=single,
  rulecolor=\color{black!25},
  xleftmargin=0.4em,
  xrightmargin=0.4em
}

\newtheorem{theorem}{Theorem}[section]
\newtheorem{proposition}[theorem]{Proposition}
\newtheorem{lemma}[theorem]{Lemma}
\newtheorem{corollary}[theorem]{Corollary}
\newtheorem{claim}[theorem]{Claim}
\theoremstyle{definition}
\newtheorem{definition}[theorem]{Definition}
\newtheorem{example}[theorem]{Example}
\newtheorem{computation}[theorem]{Exact computation}
\newtheorem{certifiedfact}{Certified finite fact}
\renewcommand{\thecertifiedfact}{C\arabic{certifiedfact}}
\crefname{certifiedfact}{Certified Fact}{Certified Facts}
\newtheorem{assumption}{Assumption}
\renewcommand{\theassumption}{A\arabic{assumption}}
\theoremstyle{remark}
\newtheorem{remark}[theorem]{Remark}
\newtheorem{externalresult}{External Result}
\renewcommand{\theexternalresult}{E\arabic{externalresult}}
\crefname{externalresult}{External Result}{External Results}

\newcommand{\Q}{\mathbb Q}
\newcommand{\Z}{\mathbb Z}
\newcommand{\N}{\mathbb N}
\newcommand{\HH}{\operatorname{HH}}
\newcommand{\KhR}{\operatorname{KhR}}
\newcommand{\SLM}{\mathcal S}
\newcommand{\im}{\operatorname{im}}
\newcommand{\id}{\operatorname{id}}
\newcommand{\Aut}{\operatorname{Aut}}
\newcommand{\Span}{\operatorname{span}}
\newcommand{\thdeg}{\operatorname{th}}
\newcommand{\qTr}{q\!\operatorname{Tr}}
\newcommand{\Sh}{\operatorname{Sh}}
\newcommand{\Status}[1]{\par\smallskip\noindent\textit{Status: #1.}\par\smallskip}
\newcommand{\ExternalStatus}{external theorem restatement}
\newcommand{\CertifiedStatus}{candidate-specific; certified by exact computation}
\newcommand{\UncertifiedStatus}{candidate-specific; not certified}
\newcommand{\eps}{\varepsilon}
\newcommand{\Hash}[1]{{\ttfamily\footnotesize\seqsplit{#1}}}

\title[An explicit degree-494 candidate class]{An explicit degree-494 candidate class for the\\
Cappell--Shaneson sphere $X(41,189,73)$}
\date{}
\subjclass[2020]{Primary 57R60, 57K18; Secondary 57R56, 57K16}
\keywords{smooth four-dimensional Poincar\'e conjecture, Cappell--Shaneson sphere,
skein lasagna module, Khovanov homology, Hochschild trace, formal verification}

% Superseded by paper/spc4-t73-candidate/main.tex (controlling manuscript).
% Retained as a historical ledger for Assumptions A1--A5 only.

\begin{document}

\begin{abstract}
We describe a candidate obstruction to the standardness of the
Cappell--Shaneson homotopy sphere associated with the matrix
\[
  A=\begin{pmatrix}0&269&1240\\0&41&189\\1&0&32\end{pmatrix}.
\]
The finite part of the calculation is exact.  A $45{,}360$-letter cabled
Artin word acts through the unreduced Burau representation on an
$88$-dimensional module over $\Z[\eps]/(\eps^7)$.  A single detector gives
\[
  D_3(v_T)=2624,\qquad D_3(\xi)=0,
\]
and a separate grading calculation places the proposed class in quantum
degree $494$.  We give all conventions, the generator matrices, the cabling
map, immutable input hashes, an independently executable reconstruction, and
the complete degree-three vector.

The passage from this finite calculation to a closed four-manifold invariant
is conditional.  We restate every external theorem with its complete
hypotheses and collect the unresolved candidate-specific geometry into five
numbered assumptions.  Under those assumptions the detector annihilates the one-, two-, and three-handle
relations, so the degree-$494$ class survives in the rational
$\mathfrak{sl}_2$ skein lasagna module.  Since the corresponding graded piece
of the standard $4$-sphere vanishes, the assumptions imply that the candidate
is not diffeomorphic to $S^4$.  The accompanying Lean development verifies
the finite arithmetic and the abstract quotient implication, but deliberately
does not instantiate the geometric assumptions.  Thus this paper is a
self-contained account of a conditional candidate proof and a precise list of
what an independent review must still establish.
\end{abstract}

\section*{Erratum to the 2 September public draft}

The 2 September public draft reported $D_3(v_T)=-59072$.  That value mixed
two endpoint coordinate systems: the vector $u$ used the THXY endpoint
indices, whereas the covector $\ell$ used the collar indices that also label
the cabled Artin word.  The two tables order the $m_2$ wickets in opposite
directions.  In the collar table, the same physical cup has endpoints $2$
and $87$, so
\[
  u=e_2-e_{87},\qquad \ell=e_{87}^*-e_2^*.
\]
With both objects in that one table, the exact calculation gives
$[\varepsilon^3]\ell(\rho(W)-I)u=-328$ and hence
$D_3(v_T)=(-2)^3(-328)=2624$.  The nonvanishing conclusion is unchanged.
Using the THXY coordinates consistently gives $-2496$; this is a coordinate
check, not a second geometric class.  The Lean literals, public receipt,
Appendix~\ref{app:vector}, and Remark~\ref{rem:serialization} have been
updated together.  The value $2624$ belongs to the currently frozen data.  If
the unresolved framing record in Assumption~\ref{ass:G1} changes the collar
word, the numerical value must be recomputed.  The superseded public PDF had
SHA-256 prefix \texttt{8E621FD9}.

\maketitle
\tableofcontents

\section{Introduction}

The smooth four-dimensional Poincar\'e conjecture asks whether every smooth
closed manifold homotopy equivalent to $S^4$ is diffeomorphic to $S^4$.  The
Cappell--Shaneson construction supplies a concrete family of homotopy
$4$-spheres from mapping tori of automorphisms of $T^3$.  Many members of this
family have been proved standard, but the construction remains a useful test
bed because the handle description is explicit enough to support exact
calculation.

This article studies the member denoted $X(41,189,73)$.  The notation refers
to a standard-form Cappell--Shaneson matrix; the actual matrix is displayed in
\cref{eq:A-matrix}.  The proposed obstruction uses the
$\mathfrak{sl}_2$ skein lasagna module of Morrison--Walker--Wedrich and its
handle-decomposition formulas.  Its logic has three steps.

\begin{enumerate}
\item Construct a class in the one-handle presentation and evaluate a divided
  cubic detector on it.  The value is the nonzero integer $2624$ and the
  class has quantum degree $494$.
\item Prove, conditional on the stated geometric identifications, that the
  same detector kills every two-handle relation and the six relations from a
  chosen system of three attaching spheres.  Linear algebra then makes the
  detector descend to the complete quotient.
\item Use the published four-handle and standard-sphere statements to compare
  the resulting degree-$494$ class with the degree-$494$ piece for $S^4$.
\end{enumerate}

The first step contains a large but elementary exact computation.  The second
step contains the principal geometric risk.  The third step is short once its
coefficient and grading hypotheses are stated correctly.  We do not blur
these three levels.  In particular, a JSON receipt or a Lean structure field
is never offered as a proof of a geometric assertion.

\subsection{Status of the main statement}

The main result of this paper is the conditional implication in
\cref{thm:main-conditional}.  It is not an unconditional disproof of the
smooth four-dimensional Poincar\'e conjecture.  The exact calculations can be
reproduced from the public repository, while the assumptions marked
``candidate-specific; not certified'' require independent geometric proof.
This division is part of the mathematical statement, not a disclaimer added
after it.

\begin{theorem}[Conditional candidate obstruction]\label{thm:main-conditional}
Assume \textup{A1--A5} below.  Then the rational
$\mathfrak{sl}_2$ skein lasagna module of $X(41,189,73)$ contains a nonzero
class in bidegree $(0,494)$.  Consequently
\[
  X(41,189,73)\not\cong_{\mathrm{diff}}S^4.
\]
Together with the Cappell--Shaneson homotopy-sphere criterion, the same
assumptions imply a counterexample to the smooth four-dimensional Poincar\'e
conjecture.
\end{theorem}

The proof appears in \cref{sec:final-descent}.  Most of the paper is devoted
to making the words ``assume A1--A5'' completely explicit.

\section{Assumption register and logical architecture}\label{sec:assumptions}

Every load-bearing statement not proved in this paper and not obtained by an
exact finite computation appears here.  The status labels have the following
fixed meanings.

\begin{itemize}
\item \emph{External theorem restatement}: a published result is used with
  no hypothesis removed and no conclusion strengthened.  The exact source
  number and scope are restated in \cref{sec:external-results}.
\item \emph{Candidate-specific; certified by exact computation}: the statement
  is about immutable finite input and can be checked by a listed program using
  exact integer or rational arithmetic.  The program does not certify any
  extra geometric interpretation.
\item \emph{Candidate-specific; not certified}: the statement identifies a
  finite model, movie, collar, or surface with the actual geometric object
  used by a published theorem.  This is a mathematical premise still needing
  proof or independent review.
\end{itemize}

\subsection{Published inputs and exact finite facts}

The proof invokes published results E1--E10, each restated with full hypotheses
in \cref{sec:external-results}: the MWW one-, two-, three-, and four-handle
formulas; the BHPW strict functoriality theorem; the BPW quantum-trace
universal property; the Horvat--Jab\l onowski attaching-sphere theorem; the
$B^4$ evaluation and smooth invariance of the skein lasagna module; and the
Cappell--Shaneson homotopy-sphere criterion.  These are citations, not new
hypotheses invented for the candidate.  The alternative van Kampen,
Wang--Mayer--Vietoris, Poincar\'e-duality, Hurewicz, and Whitehead route is
recorded separately as E11--E14.

Two finite data statements are proved by exact replay rather than assumed.

\begin{certifiedfact}[Cut and point-push input]\label{comp:finite-cut-input}
\label{ass:finite-cut}
The immutable planar-diagram input has SHA-256
\Hash{E6912A64457557469E5C691B4D57ABDBF4C45ADB05492777C574223D0C06F8A}.
Its exact traversal has $2{,}126{,}291$ crossings, $88$ cross-cut connectors,
$227$ split circle factors, and the registered oriented point-push chronology.
The reconstruction and the independent checks are given in
\cref{sec:presentation-bridge,sec:exact-computation}.
\end{certifiedfact}

\begin{certifiedfact}[Chosen-sphere finite model]\label{comp:finite-spheres}
\label{ass:chosen-sphere-data}
The three public PL data sets describe connected, oriented, genus-zero
combinatorial surface models.  Their class columns are
\[
(-1311,8608,-1)^T,\quad(-189,1241,0)^T,\quad(41,-269,1)^T,
\]
and the determinant of the resulting $3\times3$ matrix is $1$.  The row
counts, Euler calculations, and determinant are recomputed in
\cref{sec:three-handle-descent}.  The checker performs all-row PL
self-intersection and slab-gap tests only for the comb bands.  It has no
coordinates for the finger annuli, bigon, or cap, and therefore does not
certify pairwise-disjoint embedded surfaces.  That geometric statement remains
part of Assumption~\ref{ass:G4}.
\end{certifiedfact}

\subsection{The five outstanding assumptions}

\begin{assumption}[Actual presentation and legal collar]\label{ass:G1}
\label{ass:actual-presentation}\label{ass:serialization}\label{ass:collar}
\label{ass:cs-topology}
The frozen framed diagram is the actual reduced handle presentation of the
Cappell--Shaneson surgery manifold associated with \cref{eq:A-matrix}.  Its
attaching-component labels, product framings, normal fields, basepoints,
registered point-push chronology, collar, and detector insertion disk are
carried through one common framed movie.  The chronology is a legal linear
extension of that movie; coordinate changes transport the operator, vector,
shadow, covector, collar, basepoint, and insertion point simultaneously.
The finite data of C1 are the finite model of this same geometric object, and
the Cappell--Shaneson criterion E10 applies to its surgery description.
\Status{\UncertifiedStatus}
\end{assumption}

\begin{assumption}[Coefficient, endpoint, filtration, and grading binding]\label{ass:G2}
\label{ass:mww-one-handle}\label{ass:strict-functoriality}
\label{ass:quantum-trace}\label{ass:hattori-binding}
\label{ass:word-filtration}\label{ass:raw-state}
The hypotheses of E1, E8, and E9 hold for the actual cut coefficient.  The paired
annuli give the action-compatible Hattori class $\eta_R[T_1]$, whose endpoint
shadow is $u=e_2-e_{87}$ in the module of \cref{def:endpoint-module}, with
$u$ and $\ell$ read from the same collar position table.  The
registered word is the actual point-push word and lies in the third
pure-braid filtration term on the whole endpoint module.  Resealing the class
gives $s_0=(0,e_{m_2}+e_{r_{xy}})$, and the four geometric grading
contributions are $-44,+227,+315,-4$.
\Status{\UncertifiedStatus}
\end{assumption}

\begin{assumption}[All-state two-handle cubic cocone]\label{ass:G3}
\label{ass:mww-two-handle}\label{ass:two-handle-cocone}
The MWW two-handle formula E2 applies to the actual framed attaching link and
relative class.  For every finite cable state, the statewise divided-cubic
rows use one common physical-copy target, one sign convention, and the actual
core counits.  They satisfy the beta and psi equations
\cref{eq:beta-cocone,eq:psi-cocone}, and the top-cell projection is compatible
with those actual maps after the genuine quantum-trace shadow.
\Status{\UncertifiedStatus}
\end{assumption}

\begin{assumption}[Actual chosen spheres and their maps]\label{ass:G4}
\label{ass:chosen-sphere-binding}\label{ass:hj-replacement}
\label{ass:mww-three-handle}\label{ass:sphere-map-model}
The finite surfaces of C2 are embedded with the stated framings and class
coordinates in one actual post-two-handle boundary, are pairwise disjoint,
and meet every hypothesis of the Horvat--Jab\l onowski result E7.  The MWW
three-handle results E3 and E4 apply to these spheres.  On the whole source, their
actual undotted and once-dotted maps have the leading forms
$\id_{\mathrm{old}}\otimes U_j$ and
$\id_{\mathrm{old}}\otimes D_j$ after the fixed coordinate transports, and
all positive-order transport corrections enter only in order at least $h^4$
after multiplication by the cubic anomaly.
\Status{\UncertifiedStatus}
\end{assumption}

\begin{assumption}[Rational graded closure and standard-sphere comparison]\label{ass:G5}
\label{ass:four-handle}\label{ass:s4-rational}
\label{ass:diffeo-invariance}\label{ass:four-handle-grading}
The four-handle isomorphism E5 has bidegree $(0,0)$ in the convention of the
selected class.  The integral $B^4$ and $S^4$ computations E5 and E6 extend to the
rational coefficient convention so that
\[
\SLM^2_{0,q}(S^4;\Q)=0\qquad(q\ne0).
\]
The smooth-invariance isomorphism preserves that quantum grading.
\Status{\UncertifiedStatus}
\end{assumption}

\begin{center}
\begin{tabularx}{0.96\textwidth}{@{}l>{\raggedright\arraybackslash}X@{}}
\toprule
assumption & single question for independent review\\
\midrule
A1 & Is the based, framed movie---including collar and insertion point---the
actual Cappell--Shaneson presentation?\\
A2 & Does that actual coefficient produce the stated Hattori class, endpoint
vector, whole-module filtration, raw state, and grading?\\
A3 & Do the actual beta/psi maps form the all-state divided-cubic cocone?\\
A4 & Are the finite surfaces the actual attaching-sphere basis, and do their
whole-source maps have the stated leading Frobenius form?\\
A5 & Do rational coefficients and the grading convention carry the class
unchanged through the four-handle and $S^4$ comparison?\\
\bottomrule
\end{tabularx}
\end{center}
\section{Published results used, with their complete scope}\label{sec:external-results}

This section records the mathematical content imported from the literature.
The wording is a faithful restatement rather than a quotation.  In each case
the hypotheses are at least those in the cited source, and the conclusion is
not strengthened.  Candidate-specific identifications appear only in the
separate assumptions of \cref{sec:assumptions}.

\subsection{One- and two-handle formulas}

\begin{externalresult}[MWW Theorem~4.7]\label{prop:mww-4.7}
Let $k$ be a field, let $W_1=\natural^m(S^1\times B^3)$, and let
$L\subset\partial W_1$ be a null-homologous framed link.  Suppose that $L$
meets the belt sphere of the $i$th one-handle transversely in $2p_i$ points.
Cutting along these belt spheres gives a tangle $R$ in the complement of
paired balls in $S^3$.  Then the gluing map induces an isomorphism from the
direct sum, over tangles $T_i$ with the prescribed boundary, of
\[
 \KhR_N\!\left(R\cup\bigsqcup_i(T_i\sqcup\overline{T_i});k\right)
 \left\{\Bigl(\sum_i p_i\Bigr)(N-1)\right\}
\]
modulo the two gluing actions of the $3$-ball tangle categories, to
$\SLM^N_0(W_1;L;k)$.  The source states that a global grading shift is omitted
from this displayed presentation.
\end{externalresult}

\paragraph{Source and scope.}
This is External Result~E1, reproduced with the field, null-homology,
transversality, boundary, and grading-shift hypotheses that appear in MWW
Theorem~4.7 \cite[Theorem~4.7]{MWW2023Handles}.  Its candidate-specific
application is part of Assumption~\ref{ass:G2}; no such application is added
to the external statement itself.

\begin{externalresult}[MWW Definition~3.1 and Theorem~3.2]\label{prop:mww-3.2}
Let $W$ be a smooth oriented four-manifold, let $L\subset\partial W$ be a
framed link, and let $W'$ be obtained by attaching two-handles along a framed
link $K\subset\partial W$ disjoint from $L$.  For every relative class
$(\alpha,\eta)\in H_2^L(W';\Z)$, form the direct sum over $r\in\N^n$ of the
cabled summands
\[
 \SLM^N_0\!\left(W;K(r-\alpha^-,r+\alpha^+)\cup L,\eta^r\right)
 \{(1-N)(2|r|+|\alpha|)\}.
\]
Quotient by all signed-colour-preserving braid relations
$\beta_i(b)v\sim v$ and all stabilization relations
$\psi_i^{[d]}(v)\sim0$ for $d<N-1$ and
$\psi_i^{[N-1]}(v)\sim v$.  Attaching the corresponding positive and
negative parallel core disks defines an isomorphism from this cabled quotient
to $\SLM^N_0(W';L;(\alpha,\eta))$.
\end{externalresult}

\paragraph{Source and scope.}
This is exactly the scope of MWW Definition~3.1 and Theorem~3.2
\cite[Definition~3.1 and Theorem~3.2]{MWW2023Handles}.  In particular, the
anti-parallel braid relations are retained; we do not replace their braid
groups by symmetric groups.

\subsection{Three- and four-handle formulas}

\begin{externalresult}[MWW Theorem~3.7]\label{prop:mww-3.7}
Let $W$ be a smooth oriented four-manifold with boundary $Y$, let
$S\subset Y$ be an embedded $2$-sphere disjoint from a framed link
$L\subset Y$, and let $W'$ be obtained by attaching a three-handle along
$S$.  Let $J$ be an oriented equator of $S$, and let $\Delta_+$ and
$\Delta_-$ be the two hemispheres, pushed into $I\times Y$ and adjoined to
$I\times L$.  For a fixed outgoing relative class $\alpha'$, take the direct
sum over all incoming relative classes that reduce to $\alpha'$ modulo
$[S]$.  Then the three-handle map to
$\SLM^N_0(W';L';\alpha')$ is surjective, its kernel is the image of the
difference of the two hemisphere maps, and the target is their coequalizer.
\end{externalresult}

\paragraph{Source and scope.}
This is the precise content of MWW Theorem~3.7
\cite[Theorem~3.7]{MWW2023Handles}.  The direct sum over relative classes is
part of the theorem and is not suppressed here.

\begin{externalresult}[MWW Theorem~3.10]\label{prop:mww-3.10}
Suppose
$W_1\subset W_2\subset W_3\subset W_4$ is a fixed handle decomposition in
which $W_1=\natural^m(S^1\times B^3)$, $W_2$ is obtained by attaching
two-handles along $K$, $W_3$ by attaching three-handles along
$S_1,\ldots,S_p$, and $W_4$ by attaching four-handles.  Let
$L\subset\partial W_4$ be a framed link, represent the spheres by genus-zero
surfaces in $\partial W_1$ completed by two-handle cores, and fix a relative
class $\alpha'$.  Then $\SLM^N_0(W_4;L;\alpha')$ is the quotient of the
direct sum of the cabled modules over all lifts of $\alpha'$ by, for each
three-handle sphere, the undotted through $(N-2)$-dotted relations equal to
zero and the $(N-1)$-dotted relation equal to the identity.
\end{externalresult}

\paragraph{Source and scope.}
This is MWW Theorem~3.10 with its handle-decomposition, surface-realization,
link, and relative-class hypotheses \cite[Theorem~3.10]{MWW2023Handles}.
For $N=2$ it yields one undotted relation and one once-dotted-minus-identity
relation for each sphere.

\begin{externalresult}[MWW Proposition~3.4 and Corollary~3.5]\label{prop:mww-3.4}
For the empty boundary link, inclusion of a four-manifold into the result of a
three-handle attachment induces a surjection on $\SLM^N_0$, while inclusion
into the result of a four-handle attachment induces an isomorphism.  Over the
integral coefficient convention of the source,
$\SLM^N_0(S^4)\cong\Z$ and is concentrated in bidegree $(0,0)$.
\end{externalresult}

\paragraph{Source and scope.}
This is MWW Proposition~3.4 and Corollary~3.5
\cite[Proposition~3.4 and Corollary~3.5]{MWW2023Handles}.  The source
uses the two-handlebody identification of
Manolescu--Neithalath \cite[Proposition~2.1]{MN2022} in the handle formula.
The source
corollary is integral.  Its direct rational use is therefore isolated as
Assumption~\ref{ass:s4-rational} rather than being strengthened here.

\begin{externalresult}[Morrison--Walker--Wedrich Definition~5.4 and Example~5.6]\label{prop:mww-B4}
For a smooth oriented four-manifold $W$ and a framed oriented link
$L\subset\partial W$, the group $\SLM^N_0(W;L)$ is the bigraded abelian group
generated by labelled lasagna fillings modulo multilinearity, replacement of
an input ball by a filling with the same $\KhR_N$ evaluation, and isotopy
relative to the boundary.  If $W=B^4$ and $L\subset S^3=\partial B^4$, the
evaluation of fillings induces an isomorphism
\begin{equation}\label{eq:B4-evaluation}
\SLM^N_0(B^4;L)\xrightarrow{\cong}\KhR_N(S^3,L).
\end{equation}
\end{externalresult}

\paragraph{Source and scope.}
This is the invariant definition and $4$-ball evaluation in
Morrison--Walker--Wedrich \cite[Definition~5.4 and Example~5.6]{MWW2022Invariants}.
The source states the integral bigraded group.  Passage to the precise
rational convention and its compatibility with the later handle maps remains
visible in Assumption~\ref{ass:s4-rational}.

\subsection{Replacement of the three attaching spheres}

\begin{externalresult}[Horvat--Jab\l onowski Theorem~5.3]\label{prop:hj-theorem53}
Let $W$ be a smooth compact connected four-manifold with nonempty connected
boundary, equipped with a fixed handle decomposition without four-handles.
Suppose the decomposition has $k$ three-handles and
\[
 \partial W_2\cong M'\#\bigl(\#_k(S^1\times S^2)\bigr),
\]
where $M'$ is irreducible.  For pairwise-disjoint smoothly embedded
$2$-spheres $\Sigma_1,\ldots,\Sigma_k\subset\partial W_2$, the following are
equivalent: they may be isotoped, up to permutation and three--three handle
slides, to the actual attaching spheres; their exterior is connected; and
their homology classes form a $\Z$-basis of
$H_2^{\mathrm{sph}}(\partial W_2)$.  When these conditions hold, simultaneous
surgery on the spheres produces a closed $3$-manifold diffeomorphic to $M'$.
\end{externalresult}

\paragraph{Source and scope.}
This is the full statement of Theorem~5.3 in version~3 of
Horvat--Jab\l onowski \cite[Theorem~5.3]{HorvatJablonowski2026}.  The boundary
decomposition, irreducibility, number of spheres, pairwise disjointness, and
fixed-handle-decomposition hypotheses are all retained.

\subsection{Trace, functoriality, and the Cappell--Shaneson criterion}

\begin{externalresult}[BPW Proposition~3.20 and Theorem~3.21]\label{prop:bpw-trace}
For a locally pregraded endobicategory $(\mathcal C,\Sigma)$ with left duals,
the quantum horizontal trace is universal among $\Sigma$-twisted quantum
preshadows; if right duals also exist, it is a quantum shadow.  On the
$2$-category of locally pregraded endobicategories with duals and
degree-preserving structure maps, the quantum horizontal trace is a strict
$2$-functor.
\end{externalresult}

\paragraph{Source and scope.}
This combines, without enlarging either conclusion, BPW Proposition~3.20 and
Theorem~3.21 \cite[Proposition~3.20 and Theorem~3.21]{BPW2018}.

\begin{externalresult}[BHPW Theorem~4.6]\label{prop:bhpw}
For an oriented tangle $T$, the homotopy type of the Blanchet--Khovanov
bracket is an invariant of $T$ and is strictly functorial with respect to
ordinary tangle cobordisms.  Under the stated BHPW equivalence from the foam
category to the oriented Bar--Natan category, its image is isomorphic to the
ordinary Khovanov bracket.
\end{externalresult}

\paragraph{Source and scope.}
This is the strictification result in BHPW
\cite[Theorem~4.6]{BHPW2023}.  The immediately following remark in the source
labels the extension to knotted webs and four-dimensional singular foams as
conjectural.  Those objects are excluded from Assumption~\ref{ass:strict-functoriality}.

\begin{externalresult}[Iwaki Proposition~2.1]\label{prop:cs-criterion}
Let $A\in\mathrm{SL}(3,\Z)$ and let $\Sigma_A^\epsilon$ be obtained from the
mapping torus of the induced automorphism of $T^3$ by surgery on the canonical
section circle with either of the two framing parities.  Then
$\Sigma_A^\epsilon$ is a homotopy $4$-sphere if and only if
$\det(A-I)=\pm1$.
\end{externalresult}

\paragraph{Source and scope.}
This is Iwaki Proposition~2.1, which in turn cites the original
Cappell--Shaneson construction \cite[Proposition~2.1]{Iwaki2025}.  The premise
that our frozen handle diagram is this surgery manifold is not part of that
theorem; it remains Assumption~\ref{ass:actual-presentation}.

\subsection{Standard topology results represented in Lean}

The next four items are not needed once Assumption~\ref{ass:cs-topology} is
used, but they state the longer topology route represented in
\path{T73CSTopology.lean}.  They are included because a reader auditing the
Lean fields must know the exact external content behind those fields.

\begin{externalresult}[Seifert--van Kampen theorem]\label{prop:van-kampen}
Let $X=U\cup V$, where $U$, $V$, and $U\cap V$ are open and path connected,
and choose a basepoint in $U\cap V$.  Then the diagram induced by inclusion
\[
\pi_1(U\cap V)\longrightarrow\pi_1(U),
\qquad
\pi_1(U\cap V)\longrightarrow\pi_1(V)
\]
is a pushout diagram of groups: $\pi_1(X)$ is the free product of
$\pi_1(U)$ and $\pi_1(V)$ modulo the normal closure of the two images of each
element of $\pi_1(U\cap V)$.
\end{externalresult}

\paragraph{Source and scope.}
This is the standard based, path-connected form of van Kampen
\cite[Section~1.2]{Hatcher2002}.  Applying it to the Cappell--Shaneson surgery
still requires the geometric identification of the cover and of the map
$A-I$; Lean stores that application as a field.

\begin{externalresult}[Wang and Mayer--Vietoris exact sequences]\label{prop:wang-mv}
For a fibration $F\to E\to S^1$ with monodromy $f:F\to F$, there is a long
exact Wang sequence whose relevant portion is
\[
\cdots\to H_k(F)\xrightarrow{\,I-f_*\,}H_k(F)
\longrightarrow H_k(E)\longrightarrow
H_{k-1}(F)\xrightarrow{\,I-f_*\,}H_{k-1}(F)\to\cdots.
\]
For an excisive decomposition $X=A\cup B$, singular homology admits the long
exact Mayer--Vietoris sequence
\[
\cdots\to H_k(A\cap B)\to H_k(A)\oplus H_k(B)
\to H_k(X)\to H_{k-1}(A\cap B)\to\cdots,
\]
with the standard inclusion-induced maps and connecting homomorphism.
\end{externalresult}

\paragraph{Source and scope.}
These are the ordinary singular-homology exact sequences
\cite[Sections~2.2 and 4.2]{Hatcher2002}.  The candidate-specific claim that
their maps are the matrices in \cref{eq:A-minus-I,eq:H2-map} is not part of
the general theorem.

\begin{externalresult}[Integral Poincar\'e duality]\label{prop:PD}
If $M$ is a closed connected oriented topological $n$-manifold with
fundamental class $[M]\in H_n(M;\Z)$, cap product with $[M]$ gives an
isomorphism
\[
H^k(M;\Z)\xrightarrow{\cong}H_{n-k}(M;\Z)
\]
for every $k$.
\end{externalresult}

\paragraph{Source and scope.}
This is the integral oriented form of Poincar\'e duality
\cite[Section~3.3]{Hatcher2002}.  Closedness, connectedness, orientation, and
the use of integral coefficients are retained.

\begin{externalresult}[Hurewicz and Whitehead promotion]\label{prop:hurewicz-whitehead}
If a path-connected space $X$ is $(n-1)$-connected for $n\ge2$, then
$H_i(X;\Z)=0$ for $i<n$ and the Hurewicz map
$\pi_n(X)\to H_n(X;\Z)$ is an isomorphism.  If a map between simply connected
CW complexes induces an isomorphism on all integral homology groups, then it
is a homotopy equivalence.
\end{externalresult}

\paragraph{Source and scope.}
The first sentence is the Hurewicz theorem and the second is the homological
form of Whitehead's theorem for simply connected CW complexes
\cite[Sections~4.1--4.2]{Hatcher2002}.  Applied successively, they show that a
simply connected CW complex with the integral homology of $S^4$ is homotopy
equivalent to $S^4$.  The premise that the candidate has that topology is
external to the abstract theorems.

\begin{remark}[Why the short route is used]\label{rem:topology-tools}
The final proof uses Iwaki Proposition~2.1 rather than silently rebuilding the
Cappell--Shaneson calculation from these four results.  The longer route and
its Lean fields remain in \cref{tab:crosswalk}; the displayed statements make
clear that Lean does not prove their candidate-specific hypotheses.
\end{remark}

\section{The Cappell--Shaneson input}\label{sec:cs-input}

\subsection{The matrix and its elementary invariants}

\begin{definition}\label{def:A}
Let
\begin{equation}\label{eq:A-matrix}
A=
\begin{pmatrix}
0&269&1240\\
0&41&189\\
1&0&32
\end{pmatrix}
\in M_3(\Z).
\end{equation}
The symbol $X(41,189,73)$ denotes the framed surgery presentation that is
asserted in Assumption~\ref{ass:actual-presentation} to be the
Cappell--Shaneson surgery manifold $\Sigma_A^\epsilon$.  The parameter $73$
belongs to the historical standard-form naming convention; it is not an
entry of \cref{eq:A-matrix}.
\end{definition}

\begin{computation}[The two determinant tests]\label{comp:determinants}
The matrix $A$ has determinant $1$, and so does $A-I$.
\end{computation}

\begin{proof}
Expansion along the first column of $A$ gives
\begin{align}
\det A
 &=269\cdot189-1240\cdot41 \label{eq:detA-expand}\\
 &=50841-50840=1. \label{eq:detA-value}
\end{align}
Moreover
\begin{equation}\label{eq:A-minus-I}
A-I=
\begin{pmatrix}
-1&269&1240\\
0&40&189\\
1&0&31
\end{pmatrix}.
\end{equation}
Expanding along its first row gives
\begin{align}
\det(A-I)
 &=(-1)(40\cdot31)-269(0\cdot31-189\cdot1)
   +1240(0\cdot0-40\cdot1) \label{eq:detAmI-expand}\\
 &=-1240+50841-49600=1. \label{eq:detAmI-value}
\end{align}
Every operation in \cref{eq:detA-expand,eq:detAmI-expand} is integer
arithmetic.  The same identities are checked in Lean by
\texttt{T73.detA\_eq\_one} and
\texttt{T73.detAMinusI\_eq\_one}.
\end{proof}

\begin{computation}[Integral inverses]\label{comp:lattice-inverses}
The endomorphisms $A-I$ on $H_1(T^3;\Z)$ and
$(A^{-1})^T-I$ on $H_2(T^3;\Z)$ are integral automorphisms.
\end{computation}

\begin{proof}
For $B=A-I$, define
\begin{equation}\label{eq:B-inverse}
B^{-1}_{\Z}=
\begin{pmatrix}
1240&-8339&1241\\
189&-1271&189\\
-40&269&-40
\end{pmatrix}.
\end{equation}
Multiplying \cref{eq:A-minus-I} by \cref{eq:B-inverse} on either side gives
the identity matrix.  For example, the first row of $BB^{-1}_{\Z}$ is
\begin{align}
(&-1240+269\cdot189-1240\cdot40,\notag\\
 &8339-269\cdot1271+1240\cdot269,\notag\\
 &-1241+269\cdot189-1240\cdot40)=(1,0,0),
\end{align}
and the other six entries are checked by the same integer multiplication.

Since \cref{eq:detA-value} gives $A^{-1}\in M_3(\Z)$, direct inversion gives
\begin{equation}\label{eq:H2-map}
C=(A^{-1})^T-I=
\begin{pmatrix}
1311&189&-41\\
-8608&-1241&269\\
1&0&-1
\end{pmatrix}.
\end{equation}
An integral inverse is
\begin{equation}\label{eq:H2-inverse}
C^{-1}_{\Z}=
\begin{pmatrix}
-1241&-189&40\\
8339&1270&-269\\
-1241&-189&39
\end{pmatrix}.
\end{equation}
Again $CC^{-1}_{\Z}=C^{-1}_{\Z}C=I_3$.  These are the matrices used by
\path{T73CSAlgebra.lean}; the Lean proofs expand all nine entries and close
them by ring normalization.
\end{proof}

\begin{remark}[Long topology route]\label{rem:long-topology-route}
The two inverses have the expected topological consequences once the mapping
torus and surgery are identified.  Surjectivity of $A-I$ kills the natural
coinvariant presentation of the fundamental group after surgery; bijectivity
of $A-I$ and $(A^{-1})^T-I$ kills the intermediate homology in the Wang and
Mayer--Vietoris sequences.  Poincar\'e duality then gives the homology profile
of $S^4$.  Hurewicz and Whitehead promote simple connectivity plus that
homology profile to a homotopy equivalence with $S^4$.  Those implications are
represented by separate fields in \path{T73CSTopology.lean}; the field types
are shown in \cref{tab:crosswalk}.  This paper uses the shorter published
Cappell--Shaneson criterion instead, so no step of this longer route is needed
to conceal an extra geometric premise.
\end{remark}

\begin{corollary}[Conditional homotopy-sphere conclusion]\label{cor:homotopy-sphere}
Under Assumption~\ref{ass:G1} and External Result~E10,
$X(41,189,73)$ is a homotopy $4$-sphere.
\end{corollary}

\begin{proof}
By \cref{eq:detA-value}, $A\in\mathrm{SL}(3,\Z)$.  By
\cref{eq:detAmI-value}, $\det(A-I)=1$.  Assumption~\ref{ass:G1}
identifies the candidate with $\Sigma_A^\epsilon$; External Result~E10,
namely \cref{prop:cs-criterion}, then gives the conclusion.
\end{proof}

\subsection{What the topology calculation does not prove}

The determinant test proves no smooth standardness statement.  It identifies
the surgery manifold as a homotopy sphere, conditional on the presentation
binding.  The proposed obstruction below is needed precisely because the
usual homotopy invariants do not distinguish the candidate from $S^4$.
Conversely, a nonzero finite scalar attached only to the displayed matrix is
not automatically a smooth invariant.  The handle and naturality assumptions
are what connect the scalar to the closed manifold.

\section{From the mapping torus to the frozen cut tangle}\label{sec:presentation-bridge}

This section explains the geometry that Assumption~\ref{ass:actual-presentation}
must bind.  Several steps are elementary or finite; the assertion that their
composite is the actual framed Cappell--Shaneson handle presentation remains
an assumption.

\subsection{The mapping-torus relations}

Let $x_1,x_2,x_3$ be the coordinate generators of
$\pi_1(T^3)=\Z^3$ and let $t$ denote the base circle of the mapping torus.
The mapping-torus presentation has relations
\begin{equation}\label{eq:mapping-torus-relations}
m_i=t\,\phi_A(x_i)\,t^{-1}x_i^{-1},\qquad i=1,2,3,
\end{equation}
where $\phi_A$ is induced by the column action of $A$.  Aitchison--Rubinstein
describe each $m_i$ as the boundary of a product annulus assembled from a
bottom coordinate arc, its top image, and two vertical connecting arcs
\cite[pp.~5--9]{AitchisonRubinstein1984}.  The product annulus, rather than a
blackboard writhe, specifies the framing used here.

The first column of \cref{eq:A-matrix} is $e_3$, so
$\phi_A(x_1)=x_3$.  After the section/base-handle cancellation,
\cref{eq:mapping-torus-relations} gives
\begin{equation}\label{eq:m1-after-first-cancel}
m_1=x_3x_1^{-1}.
\end{equation}
Cancelling this complementary one--two handle pair identifies $x_1$ with
$x_3$.  Under that substitution the three commutator components become
\begin{equation}\label{eq:commutator-substitution}
[x_1,x_2]\mapsto[x_3,x_2],\qquad
[x_2,x_3]\mapsto[x_2,x_3],\qquad
[x_3,x_1]\mapsto[x_3,x_3]=1.
\end{equation}
The last component bounds the two nested product bigons obtained by cancelling
$x_3x_3x_3^{-1}x_3^{-1}$; it is a split zero-framed component in the product
model.

\begin{computation}[Reduced word data]\label{comp:reduced-words}
The frozen integer word builder applies \cref{eq:mapping-torus-relations} and
the two substitutions above to obtain a reduced word $m_2$ of length $311$
with exponent sum $(40,269)$ in $(x_2,x_3)$, and a reduced word $m_3$ of
length $1460$ with exponent sum $(189,1271)$.
\end{computation}

\begin{proof}
The words are generated from the columns $Ae_2=(269,41,0)^T$ and
$Ae_3=(1240,189,32)^T$ by the Christoffel-path convention in the public
builder, followed by free reduction after $x_1=x_3$.  The checker
\path{scripts/verify_public_geometry_evidence.py} regenerates the paths and
their exponent sums from \path{evidence/public_geometry/cs_presentation.py}.
Because the output is an integer word, both the length and the exponent sums
are exact.  The statement that these paths are the transported top arcs of
the actual framed product annuli is not a consequence of word equality and
remains in Assumption~\ref{ass:actual-presentation}.
\end{proof}

\subsection{A relative descending criterion}

\begin{lemma}[Relative global descending lemma]\label{lem:global-descending}
Let $T$ be a marked tangle in a $3$-ball.  Suppose its boundary endpoints and
pairing are fixed; its components carry a total order; on each component the
branch encountered earlier from the marked initial point is over at every
self-crossing; and at every mutual crossing an earlier component is over a
later one.  Then $T$ is ambient-isotopic relative to the boundary to the
boundary-parallel trivial tangle with the same marked endpoint pairing.
\end{lemma}

\begin{proof}
Place the first component in a thin top layer.  Traverse it from the marked
initial endpoint.  At a self-crossing, the already traversed branch lies over
the untraversed branch; at a mutual crossing, the whole component lies over
every later component.  The traversed arc can therefore be pushed into a
boundary collar without crossing the remaining tangle.  Continuing produces
a bridge disk disjoint from all later components.  Remove the first component
and repeat.  Induction gives disjoint bridge disks for every component, hence
an ambient isotopy relative to the boundary.  The fixed endpoint pairing
determines the unique marked boundary-parallel target.
\end{proof}

\begin{computation}[Descending check for the frozen PD]\label{comp:descending}
The exact traversal of the frozen $2{,}126{,}291$-crossing PD input satisfies
the hypotheses of Lemma~\ref{lem:global-descending}.  The self-crossing counts
for $(r_{xy},r_{yz},m_2,m_3,r_{zx})$ are
\begin{equation}\label{eq:self-cross-counts}
(9,3,64631,1445582,0),
\end{equation}
and every mutual crossing follows the strict component order
\begin{equation}\label{eq:height-order}
r_{xy}>r_{yz}>m_2>m_3.
\end{equation}
\end{computation}

\begin{proof}
The script \path{scripts/verify_t73_global_descending.py} reads every crossing,
reconstructs the cyclic traversal including the closing connector, checks the
over/under condition at every self-crossing, and checks
\cref{eq:height-order} at every mutual crossing.  It recomputes the counts in
\cref{eq:self-cross-counts} rather than trusting summary fields.  The input
identity is the SHA-256 in Certified Fact~\ref{comp:finite-cut-input}.
\end{proof}

The descending lemma proves that the frozen cut tangle is trivial relative to
its marked endpoints.  It does not prove that the tangle is the cut of the
actual Aitchison--Rubinstein product presentation.  Assumption
\ref{ass:actual-presentation} supplies that last whole-link, framing-preserving
identification.  Laudenbach--Poenaru's theorem permits a boundary
self-diffeomorphism of $\#_4(S^1\times S^2)$ to extend over the four-dimensional
one-handlebody \cite{LaudenbachPoenaru1972}; it does not identify the finite
diagram without the preceding geometric construction.

\section{Coefficient traces and the selected Hattori class}\label{sec:hattori}

\subsection{Coefficient Hochschild homology}

\begin{definition}[Linear category and bimodule]\label{def:category-bimodule}
Let $\mathcal C$ be a small $\Q$-linear category.  A coefficient bimodule
$M$ assigns a vector space $M(T,T')$ to each ordered pair of objects and has
commuting left and right actions
\begin{equation}\label{eq:bimodule-actions}
\mathcal C(S,T)\otimes M(T,T')\longrightarrow M(S,T'),
\qquad
M(T,T')\otimes\mathcal C(T',U)\longrightarrow M(T,U).
\end{equation}
Associativity and unit conditions are the usual ones.
\end{definition}

\begin{definition}[Coefficient trace]\label{def:coefficient-HH0}
The coefficient zeroth Hochschild homology is
\begin{equation}\label{eq:HH0}
\HH_0(\mathcal C;M)=
\left(\bigoplus_{T\in\operatorname{Ob}\mathcal C}M(T,T)\right)
\Big/
\Span_{\Q}\{f\cdot m-m\cdot f\},
\end{equation}
where $f:S\to T$ and $m\in M(T,S)$ range over every cyclically composable
pair.  If $m\in M(T,T)$ is a diagonal cycle, its image in
\cref{eq:HH0} is denoted $[m]_T$.  When $m$ is the identity-type element
obtained from a representable coefficient, we call $[m]_T$ its
\emph{Hattori class}.
\end{definition}

This definition is deliberately concrete.  It is the quotient implemented in
\path{Smooth4PC/CoefficientTrace.lean}.  The quotient universal property is
ordinary linear algebra: a linear map on the direct sum descends if and only
if it kills every vector $f\cdot m-m\cdot f$.

\begin{lemma}[Descent through the coefficient trace]\label{lem:HH0-descent}
Let $r_T:M(T,T)\to V$ be a family of linear maps satisfying
\begin{equation}\label{eq:cyclicity}
r_S(f\cdot m)=r_T(m\cdot f)
\end{equation}
for all cyclically composable $f$ and $m$.  Then there is a unique linear map
$\bar r:\HH_0(\mathcal C;M)\to V$ whose restriction to $M(T,T)$ is $r_T$.
\end{lemma}

\begin{proof}
Let $r$ be the sum of the $r_T$ on the algebraic direct sum.  By
\cref{eq:cyclicity},
$r(f\cdot m-m\cdot f)=0$ for every generator of the relation space in
\cref{eq:HH0}.  Hence that relation space lies in $\ker r$, so $r$ factors
uniquely through the quotient.
\end{proof}

\subsection{The Frobenius factor and the diagonal element}

\begin{definition}[Rank-two Frobenius algebra]\label{def:frobenius}
Throughout the finite $\mathfrak{sl}_2$ calculation we use
\begin{equation}\label{eq:frob-algebra}
\mathsf A=\Q[X]/(X^2),\qquad
\epsilon(1)=0,\qquad \epsilon(X)=1,
\end{equation}
with coproduct
\begin{equation}\label{eq:frob-coproduct}
\Delta(1)=1\otimes X+X\otimes1,\qquad
\Delta(X)=X\otimes X.
\end{equation}
The $227$ labels below mean the tensor
$X^{\otimes227}$, not the product $X^{227}=0$ in $\mathsf A$.
\end{definition}

\begin{definition}[Selected object and class]\label{def:vT}
Let $U:P_{86}\to P_{88}$ be the physical oriented one-cup tangle.  Let
$B_{\mathrm{act}}\in\Aut_{\mathcal C}(P_{88})$ be the automorphism in
Assumption~\ref{ass:hattori-binding}, and set
\begin{equation}\label{eq:T1}
T_1=B_{\mathrm{act}}^{-1}U.
\end{equation}
The balanced coefficient equivalence is denoted
\begin{equation}\label{eq:H-balanced}
H_{T,T'}:M_R(T,T')\xrightarrow{\cong}
\operatorname{Hom}_{\mathcal C}(B_{\mathrm{act}}T,B_{\mathrm{act}}T')
\{-44\}\otimes\mathsf A^{\otimes227}.
\end{equation}
Define
\begin{equation}\label{eq:vT}
v_T=H_{T_1,T_1}^{-1}
\bigl(\id_U\otimes X^{\otimes227}\bigr),
\qquad
\eta_R[T_1]=[v_T]_{T_1}\in\HH_0(\mathcal C;M_R).
\end{equation}
\end{definition}

Equation~\eqref{eq:vT} defines the algebraic element once
\cref{eq:H-balanced} is given.  The substantive geometric statement is that
\cref{eq:H-balanced} is the coefficient equivalence arising from the actual
paired-annulus cut, with compatible left and right actions.  That statement is
Assumption~\ref{ass:hattori-binding}.

For comparison with the earlier relative input, factor
$B_{\mathrm{act}}=WF_\Omega$ and define
\begin{equation}\label{eq:T0T1-factor}
T_0=F_\Omega^{-1}U,\qquad
T_1=F_\Omega^{-1}W^{-1}U.
\end{equation}
Then associativity, without any commutativity assumption, gives
\begin{equation}\label{eq:BT0BT1}
B_{\mathrm{act}}T_0=WU,\qquad B_{\mathrm{act}}T_1=U.
\end{equation}
The first equality follows from
$WF_\Omega F_\Omega^{-1}U=WU$; the second follows from
$WF_\Omega F_\Omega^{-1}W^{-1}U=U$.

\begin{definition}[Relative class]\label{def:xi}
Let $s_{\mathrm{inv}}\in\{\pm1\}$ be the strict functorial sign assigned to
the fixed cancellation movie for $WW^{-1}$.  The relative class is
\begin{equation}\label{eq:xi}
\xi=\eta_R[T_0]-s_{\mathrm{inv}}\eta_R[T_1].
\end{equation}
It is not the selected input $\eta_R[T_1]$.
\end{definition}

\section{The endpoint module and the Burau action}\label{sec:endpoint}

\subsection{The \texorpdfstring{$88$}{88}-dimensional module}

\begin{definition}[Endpoint module]\label{def:endpoint-module}
Let
\begin{equation}\label{eq:R-eps}
R_\eps=\Z[\eps]/(\eps^7),\qquad t=1+\eps,
\end{equation}
and let
\begin{equation}\label{eq:E88}
E_\eps=R_\eps^{88}=\bigoplus_{i=0}^{87}R_\eps e_i.
\end{equation}
One may identify this free module with the weight-$86$ subspace of the
$88$-fold tensor power of the two-dimensional fundamental module: $e_i$
corresponds to the tensor having the lower-weight basis vector in position
$i$ and the upper-weight basis vector in every other position.  This explains
the dimension $\binom{88}{1}=88$.  The identification of this algebraic model
with the actual endpoint shadow is part of
Assumption~\ref{ass:hattori-binding}.
\end{definition}

\begin{definition}[Unreduced Burau representation]\label{def:burau}
For the Artin generator $\sigma_i$ of $B_{88}$, with $1\le i\le87$, define
$\rho_t(\sigma_i)$ to be the identity off the span of
$e_{i-1},e_i$ and to have the block
\begin{equation}\label{eq:burau-positive}
\rho_t(\sigma_i)\big|_{\langle e_{i-1},e_i\rangle}
=
\begin{pmatrix}1-t&t\\1&0\end{pmatrix}
=
\begin{pmatrix}-\eps&1+\eps\\1&0\end{pmatrix}.
\end{equation}
Its inverse has block
\begin{equation}\label{eq:burau-negative}
\rho_t(\sigma_i^{-1})\big|_{\langle e_{i-1},e_i\rangle}
=
\begin{pmatrix}0&1\\t^{-1}&1-t^{-1}\end{pmatrix},
\end{equation}
where
\begin{equation}\label{eq:t-inverse}
t^{-1}=1-\eps+\eps^2-\eps^3+\eps^4-\eps^5+\eps^6
\quad\text{in }R_\eps.
\end{equation}
For a word $w=\sigma_{i_1}^{s_1}\cdots\sigma_{i_m}^{s_m}$ we use the left
action
\begin{equation}\label{eq:word-action-order}
\rho_t(w)v=\rho_t(\sigma_{i_1}^{s_1})\cdots
\rho_t(\sigma_{i_m}^{s_m})v.
\end{equation}
Thus an implementation acting in place on a column vector must traverse the
stored letters from right to left.
\end{definition}

\begin{lemma}\label{lem:burau-inverse}
The blocks in \cref{eq:burau-positive,eq:burau-negative} are mutual inverses
over $R_\eps$.
\end{lemma}

\begin{proof}
Direct multiplication gives
\begin{equation}\label{eq:burau-inverse-check}
\begin{pmatrix}1-t&t\\1&0\end{pmatrix}
\begin{pmatrix}0&1\\t^{-1}&1-t^{-1}\end{pmatrix}
=
\begin{pmatrix}1&(1-t)+t-1\\0&1\end{pmatrix}=I_2.
\end{equation}
The reverse product is the same calculation.  Formula~\eqref{eq:t-inverse}
is the geometric-series inverse of $1+\eps$ modulo $\eps^7$.
\end{proof}

\subsection{From primitive crossings to the cabled word}

The public input stores $252$ primitive crossing rows, not the expanded
$11{,}340$- or $45{,}360$-letter words.  Each row records a sign
$s\in\{\pm1\}$, two paired-strand labels $1\le i<j\le44$, and one of two geometries.

\begin{definition}[Row expansion]\label{def:row-expansion}
For $i<j$ define the two signed pure-row words
\begin{align}
L_{ij}^{s}
 &=\sigma_{j-1}\sigma_{j-2}\cdots\sigma_{i+1}
   \sigma_i^{s}\sigma_i^{s}
   \sigma_{i+1}^{-1}\cdots\sigma_{j-2}^{-1}\sigma_{j-1}^{-1},
   \label{eq:L-row}\\
R_{ij}^{s}
 &=\sigma_i\sigma_{i+1}\cdots\sigma_{j-2}
   \sigma_{j-1}^{s}\sigma_{j-1}^{s}
   \sigma_{j-2}^{-1}\cdots\sigma_{i+1}^{-1}\sigma_i^{-1}.
   \label{eq:R-row}
\end{align}
The $44$-strand word $B_{44}$ is obtained by taking the $R$ rows on segment
$6$ in increasing horizontal coordinate, followed by the $L$ rows on segment
$2$ in decreasing horizontal coordinate, and replacing each row by
\cref{eq:L-row,eq:R-row}.
\end{definition}

\begin{definition}[Two-cabling]\label{def:cabling}
The cabling homomorphism on letters is
\begin{equation}\label{eq:cabling}
c(\sigma_i)=\sigma_{2i}\sigma_{2i+1}\sigma_{2i-1}\sigma_{2i},
\qquad
c(\sigma_i^{-1})=c(\sigma_i)^{-1}.
\end{equation}
The $88$-strand word is $B_{88}=c(B_{44})$.
\end{definition}

The use of ``homomorphism'' in \cref{def:cabling} is algebraic: the code
applies \cref{eq:cabling} letter by letter and reverses and changes signs for
an inverse letter.  The geometric statement that this is the actual oriented
cabling of the point-push movie belongs to
Assumptions~\ref{ass:word-filtration} and \ref{ass:serialization}.

\begin{computation}[Word identities]\label{comp:word-identities}
The exact reconstruction gives
\begin{equation}\label{eq:word-lengths}
|B_{44}|=11340,\qquad |B_{88}|=45360.
\end{equation}
The canonical JSON SHA-256 values are, respectively,
\begin{quote}
\Hash{7C2D2F792C2672221A76CAF08A71F560AF0CB7654B4D537C00FAD00B16EFA187}
\end{quote}
and
\begin{quote}
\Hash{C13D0A3BA9B05F41A6D2C5B4AB12DDECDABEFC1883835891AD3BD2B8955B5FFB}.
\end{quote}
The first forty letters are
\begin{align}
B_{44}[0{:}40]&=(1,2,3,4,5,6,7,8,9,10,11,12,13,14,15,16,17,18,19,20,
\notag\\[-2pt]
&\qquad21,22,23,24,25,26,27,28,29,30,31,32,33,34,35,36,37,38,39,40),
\label{eq:B44-first40}\\
B_{88}[0{:}40]&=(2,3,1,2,4,5,3,4,6,7,5,6,8,9,7,8,10,11,9,10,
\notag\\[-2pt]
&\qquad12,13,11,12,14,15,13,14,16,17,15,16,18,19,17,18,20,21,19,20).
\label{eq:B88-first40}
\end{align}
\end{computation}

\begin{proof}
The reconstruction is precisely Definitions~\ref{def:row-expansion} and
\ref{def:cabling}.  Canonical JSON means UTF-8 encoding with sorted keys where
applicable and list separators \texttt{,} and \texttt{:} without spaces.  The
command in \cref{sec:reproduction} recomputes the two lengths, the two hashes,
and the displayed prefixes from the primitive rows.  Neither expanded word is
stored in the input.
\end{proof}

\section{The detector and the exact cubic}\label{sec:exact-computation}

\subsection{One definition for both inputs}

Let $R_h=\Q[h]/(h^7)$ and substitute
\begin{equation}\label{eq:parameter-change}
q=1+h,\qquad t=q^{-2},\qquad \eps=t-1.
\end{equation}
Let $\Sh_h$ denote the quantum trace shadow followed by the fixed endpoint
identification, let $P_{86}$ be the rational projector onto the through-$86$
top cell, and let $\widehat C(h)$ be the normalized endpoint cap row.  These
maps are typed on the full endpoint target by Assumptions
\ref{ass:hattori-binding} and \ref{ass:two-handle-cocone}.

\begin{definition}[The detector]\label{def:detector}
For a quantum trace class $x$, define the single series
\begin{equation}\label{eq:D-h}
\mathcal D_h(x)=
\widehat C(h)\bigl(\rho_h(W)-I\bigr)P_{86}\Sh_h(x).
\end{equation}
Under Assumption~\ref{ass:word-filtration}, the image is contained in
$h^3R_h$.  The divided cubic detector on the ordinary coefficient trace is
\begin{equation}\label{eq:D3}
D_3(\bar x)=\left.\frac{\mathcal D_h(x)}{h^3}\right|_{h=0},
\end{equation}
where $x$ is any quantum lift of $\bar x$.
\end{definition}

\begin{lemma}[Lift independence]\label{lem:lift-independence}
Formula~\eqref{eq:D3} is independent of the quantum lift.
\end{lemma}

\begin{proof}
If $x'$ and $x$ have the same specialization at $h=0$, then
$x'-x=hy$.  Write $\mathcal D_h=h^3\mathcal D'_h$.  Linearity gives
\begin{equation}\label{eq:lift-difference}
\mathcal D'_h(x')-\mathcal D'_h(x)=h\mathcal D'_h(y),
\end{equation}
whose reduction modulo $h$ is zero.  Thus \cref{eq:D3} has the same value on
$x$ and $x'$.
\end{proof}

\begin{lemma}[Specialization of the quantum trace]\label{lem:qtrace-specialization}
Let $\qTr_q(\mathcal C;M)$ be the quotient of
$\bigoplus_T M(T,T)\otimes\Q[q,q^{-1}]$ by
\begin{equation}\label{eq:qtrace-relation}
L_f(m)-q^{|f|}R_f(m).
\end{equation}
Then
\begin{equation}\label{eq:qtrace-to-HH0}
\qTr_q(\mathcal C;M)\otimes_{\Q[q,q^{-1}]}\Q[q,q^{-1}]/(q-1)
\cong\HH_0(\mathcal C;M).
\end{equation}
\end{lemma}

\begin{proof}
Tensor product is right exact.  Applying it to the cokernel presentation of
$\qTr_q$ replaces $q^{|f|}$ by $1$, so every generator
\cref{eq:qtrace-relation} becomes $L_f(m)-R_f(m)$.  The resulting cokernel is
exactly \cref{eq:HH0}.  No flatness assumption on $\qTr_q$ is needed.
\end{proof}

Under the endpoint normalization of the public input, the composite row
$\widehat C(0)P_{86}$ is
\begin{equation}\label{eq:ell}
\ell=e_{87}^{*}-e_2^{*},
\end{equation}
and the selected shadow vector is
\begin{equation}\label{eq:u}
u=e_2-e_{87}.
\end{equation}
Here $e_i^*(e_j)=\delta_{ij}$.  Both formulas use the collar-coordinate table
\path{data/B88_POSITION_TO_PASSAGE_TABLE.json}; its SHA-256 is
\Hash{119C7E9E74AE6C820DA72A84CDFD5D445D81E6C3AACCC209C25D37C323961508}.
In that table positions $2$ and $87$ are the negative copy of wicket $W_2$
and the positive copy of wicket $W_{44}$, respectively.

\subsection{The exact vector calculation}

\begin{computation}[Burau difference on the selected vector]\label{comp:burau-vector}
In $E_\eps$, the first nonzero coefficient of
$(\rho_t(B_{88})-I)u$ occurs in degree $3$.  Its full coefficient vector is
\begin{equation}\label{eq:eps3-vector-short}
[\eps^3](\rho_t(B_{88})-I)u=(a_0,\ldots,a_{87})^T,
\end{equation}
where the $88$ integers $a_i$ are printed in
\cref{app:vector}.  Contracting with \cref{eq:ell} gives
\begin{equation}\label{eq:ell-eps-series}
\ell(\rho_t(B_{88})-I)u
=-328\eps^3+14596\eps^4-410246\eps^5+9595271\eps^6.
\end{equation}
\end{computation}

\begin{proof}
Start with the sparse column $u$ in \cref{eq:u}.  Traverse the letters of
$B_{88}$ from right to left, applying \cref{eq:burau-positive} for a positive
letter and \cref{eq:burau-negative} for a negative letter.  All polynomial
products are reduced modulo $\eps^7$ after each operation.  Subtract the
original column.  Finally take coordinate $87$ minus coordinate $2$ as in
\cref{eq:ell}.  This deterministic algorithm is implemented in
\texttt{scripts/recompute\_t73\_delta3.py}; it uses Python integers only.
The input SHA-256 is
\begin{quote}
\Hash{0F0F4A568D966367313C29C25CEDDE4FD9ABAD32B9375519E9CDF9168AE0A764}.
\end{quote}
The program rejects a mismatch in the primitive-row hash, either word hash,
the representation block, the cabling rule, or the Hattori word binding.
Thus \cref{eq:ell-eps-series} is independently reimplementable from the
formulas in this paper and the immutable primitive rows.
\end{proof}

\begin{lemma}[Parameter expansion]\label{lem:eps-of-h}
In $\Z[h]/(h^7)$,
\begin{equation}\label{eq:eps-h-expansion}
\eps=(1+h)^{-2}-1=-2h+3h^2-4h^3+5h^4-6h^5+7h^6.
\end{equation}
\end{lemma}

\begin{proof}
The binomial series gives
$(1+h)^{-2}=\sum_{k\ge0}(-1)^k(k+1)h^k$.
Subtracting $1$ and truncating after degree $6$ yields
\cref{eq:eps-h-expansion}.
\end{proof}

\begin{lemma}[Commutator order]\label{lem:commutator-order}
Let $R$ be an $h$-adically filtered ring.  If
$P=I+O(h^a)$ and $Q=I+O(h^b)$ are invertible matrices, then
\begin{equation}\label{eq:commutator-order}
PQP^{-1}Q^{-1}=I+O(h^{a+b}).
\end{equation}
Consequently, if every generator of a pure-braid subgroup acts as
$I+O(h)$, then its $r$th lower-central subgroup acts as $I+O(h^r)$.
\end{lemma}

\begin{proof}
Write $P=I+h^aP_0$ and $Q=I+h^bQ_0$.  Modulo $h^{a+b}$, the inverse series
are $P^{-1}=I-h^aP_0+O(h^{2a})$ and
$Q^{-1}=I-h^bQ_0+O(h^{2b})$.  Multiplication shows that all terms of order
strictly below $a+b$ cancel in $PQP^{-1}Q^{-1}$.  The lower-central claim
follows by induction from
$\Gamma_{r+1}=[\Gamma_r,\Gamma_1]$.
\end{proof}

Lemma~\ref{lem:commutator-order} proves the filtration implication.  It does
not prove that the geometric point-push word belongs to $\Gamma_3$; that
membership and the identification of the word with the actual movie are the
unproved parts of Assumption~\ref{ass:word-filtration}.

\begin{proposition}[Selected value]\label{prop:D3-vT}
Under Assumption~\ref{ass:hattori-binding}, the detector of
\cref{def:detector} satisfies
\begin{equation}\label{eq:D3-vT}
D_3(\eta_R[T_1])=2624\ne0.
\end{equation}
\end{proposition}

\begin{proof}
Assumption~\ref{ass:hattori-binding} and \cref{eq:BT0BT1} identify the
leading endpoint shadow of $\eta_R[T_1]$ with $u$.  By
\cref{eq:ell-eps-series}, the coefficients below $\eps^3$ vanish.  Hence only
the cubic term contributes to the coefficient of $h^3$.  By
\cref{eq:eps-h-expansion}, $[h]\eps=-2$, so
\begin{equation}\label{eq:D3-vT-calc}
[h^3]\,(-328)\eps^3=(-328)(-2)^3=2624.
\end{equation}
The last integer is nonzero.
\end{proof}

\begin{computation}[The squared shell]\label{comp:squared-shell}
Applying the same operator a second time gives
\begin{equation}\label{eq:squared-series}
\ell(\rho_t(B_{88})-I)^2u=-102729600\eps^6
\quad\text{in }R_\eps.
\end{equation}
\end{computation}

\begin{proof}
Use the output column of \cref{comp:burau-vector} as the input to the same
right-to-left word action, subtract that input column, and contract with the
same covector \cref{eq:ell}.  This is the second calculation performed by the
public script.  No second detector definition is introduced.
\end{proof}

\begin{computation}[The complete scalar profiles in the $h$ coordinate]\label{comp:h-profiles}
Substitution of \cref{eq:eps-h-expansion} into
\cref{eq:ell-eps-series,eq:squared-series} yields, through degree six,
\begin{align}
\ell(\rho_h(W)-I)u
={}&2624h^3+221728h^4+11760112h^5
\notag\\
&+520583560h^6, \label{eq:eta-h-profile}\\
\ell(\rho_h(W)-I)^2u
={}&-6574694400h^6. \label{eq:xi-h-profile}
\end{align}
\end{computation}

\begin{proof}
Polynomial substitution in the quotient $\Z[h]/(h^7)$ is finite.  For
example, the cubic coefficient of the first line is
$(-328)(-2)^3=2624$.  For the second line,
$\eps^6=(-2)^6h^6+O(h^7)=64h^6+O(h^7)$, and
$-102729600\cdot64=-6574694400$.  Expanding the remaining powers and collecting
coefficients gives \cref{eq:eta-h-profile}; the public script performs the
same convolution with integer lists.
\end{proof}

\begin{proposition}[The relative input is killed]\label{prop:D3-xi}
With $\xi$ defined by \cref{eq:xi} and with the same detector
\cref{eq:D3},
\begin{equation}\label{eq:D3-xi}
D_3(\xi)=0.
\end{equation}
\end{proposition}

\begin{proof}
The cancellation sign in \cref{def:xi} is chosen so that
\begin{equation}\label{eq:xi-shadow}
P_{86}\Sh_h(\xi)=(\rho_h(W)-I)u.
\end{equation}
Substitution of \cref{eq:xi-shadow} into the single definition
\cref{eq:D-h} produces the square
$\ell(\rho_h(W)-I)^2u$.  By \cref{eq:squared-series} it begins in order
$\eps^6$, hence in order $h^6$.  Its $h^3$ coefficient is zero.
\end{proof}

\begin{remark}[Why the two values are not contradictory]\label{rem:vT-xi}
The plain shadow coefficient $[h^3]\ell\Sh_h(\xi)$ equals $2624$, but that
plain shadow is not $D_3(\xi)$.  The detector already contains the outer
factor $\rho_h(W)-I$.  Thus the selected class $\eta_R[T_1]$ receives one
factor and has value $2624$, whereas $\xi$ already contains one factor and
the detector supplies a second, yielding zero in degree three.
\end{remark}

\section{The grading calculation}\label{sec:grading}

\begin{proposition}[Quantum degree]\label{prop:degree-494}
Under Assumption~\ref{ass:raw-state}, the selected raw class has quantum
degree $494$.
\end{proposition}

\begin{proof}
There are four contributions.  First, the $44$ paired boundary passages give
the built-in Hom normalization $44(N-1)=44$ for $N=2$; the raw closure removes
it, contributing
\begin{equation}\label{eq:degree-minus44}
d_1=-44.
\end{equation}
Second, the $227$ separately labelled factors $X^{\otimes227}$ each contribute
$+1$ in the adopted convention:
\begin{equation}\label{eq:degree-plus227}
d_2=227.
\end{equation}
Third, the one-handle theorem contributes the sum of the $44$ $y$-pairs and
$271=44+227$ $z$-pairs:
\begin{equation}\label{eq:degree-plus315}
d_3=44+271=315.
\end{equation}
Finally the selected cabled state has $\alpha=0$, $|r|=2$, and $N=2$, so the
MWW cabled shift is
\begin{equation}\label{eq:degree-minus4}
d_4=(1-N)(2|r|+|\alpha|)=(-1)\cdot4=-4.
\end{equation}
Adding \cref{eq:degree-minus44,eq:degree-plus227,eq:degree-plus315,eq:degree-minus4}
gives
\begin{equation}\label{eq:degree-total}
d_1+d_2+d_3+d_4=-44+227+315-4=494.
\end{equation}
\end{proof}

The integer arithmetic in \cref{eq:degree-total} is checked in Lean.  Lean
does not prove that these four integers are the correct geometric grading
contributions for the actual filling; that binding is exactly why
Assumption~\ref{ass:raw-state} remains visible.

\section{Descent through the one- and two-handle relations}\label{sec:two-handle-descent}

\subsection{The through-degree separation}

Let $F_d$ denote the span of endpoint diagrams of through degree at most $d$.
Composition on either side cannot increase through degree.  The selected cup
$U$ has through degree $86$.  The maps $\psi_i^{[0]}$ and
$\psi_i^{[1]}$ also preserve through degree $86$ and add the same one-cup
factor.  Thus neither through degree nor cup count distinguishes their images.
Their statewise distinction in Proposition~\ref{prop:beta-psi} comes from the
core counit of Lemma~\ref{lem:counit-identities}.  The following elementary
filtration lemma controls only relation terms that are independently known to
lie in the low-cell subspace $F_{84}$.

\begin{lemma}[Filtration argument]\label{lem:through-filtration}
Suppose $P_{86}$ is a linear projector with image the through-$86$ top cell
and kernel containing $F_{84}$.  If $I\subset F_{84}$ is closed under left and
right endpoint composition, then $P_{86}(I)=0$.  In particular, no vector in
$I$ can change $D_3(v_T)$.
\end{lemma}

\begin{proof}
The inclusion $I\subset F_{84}\subset\ker P_{86}$ gives $P_{86}(I)=0$.
The last statement follows by linearity of every map after $P_{86}$ in
\cref{eq:D-h}.
\end{proof}

The point is not that a physical cap itself kills every lower cell.  It need
not.  The projection is a linear postcomposition after the genuine quantum
trace shadow.  Its compatibility with the actual cabled relations is part of
Assumption~\ref{ass:two-handle-cocone}.

\subsection{The local Frobenius calculation}

For $b\ge1$, let $\Delta^{(b-1)}:\mathsf A\to\mathsf A^{\otimes b}$ be the
iterated coproduct, with $\Delta^{(0)}=\id$.  Put
\begin{align}
D_b&=\Delta^{(b-1)}(X)=X^{\otimes b}, \label{eq:Db}\\
U_b&=\Delta^{(b-1)}(1)
=\sum_{a=0}^{b-1}X^{\otimes a}\otimes1\otimes
X^{\otimes(b-1-a)}. \label{eq:Ub}
\end{align}

\begin{lemma}[Counit identities]\label{lem:counit-identities}
For every $b\ge1$,
\begin{equation}\label{eq:counit-identities}
\epsilon^{\otimes b}(U_b)=0,
\qquad
\epsilon^{\otimes b}(D_b)=1.
\end{equation}
\end{lemma}

\begin{proof}
Formula~\eqref{eq:frob-coproduct} implies
$\Delta^{(b-1)}(X)=X^{\otimes b}$ by induction.  Applying
$\epsilon^{\otimes b}$ gives $1^b=1$.  The corresponding induction for $1$
gives the sum in \cref{eq:Ub}; every summand contains exactly one factor
$1$, and $\epsilon(1)=0$.  Hence every summand evaluates to zero.
\end{proof}

\subsection{The cabled relation space}

Let $\mathcal C_s$ be the raw summand at a finite cable state $s$ and set
$\mathcal C_{\mathrm{raw}}=\bigoplus_s\mathcal C_s$.  The two-handle relation
space $\mathcal R_{2h}$ is spanned by
\begin{equation}\label{eq:two-handle-generators}
\beta_i(b)x-x,\qquad \psi_i^{[0]}x,\qquad
\psi_i^{[1]}x-x
\end{equation}
for every attaching component, state, admissible signed-colour-preserving braid, and source
vector.

\begin{proposition}[Cubic beta and psi equations]\label{prop:beta-psi}
Under Assumptions~\ref{ass:word-filtration} and
\ref{ass:two-handle-cocone}, the statewise rows $\Lambda_s^{(3)}$ satisfy
\begin{align}
\Lambda_s^{(3)}\beta_i(b)&=\Lambda_s^{(3)},
\label{eq:beta-cocone}\\
\Lambda_{s+e_i}^{(3)}\psi_i^{[0]}&=0,
\qquad
\Lambda_{s+e_i}^{(3)}\psi_i^{[1]}=\Lambda_s^{(3)}.
\label{eq:psi-cocone}
\end{align}
Consequently $\mathcal R_{2h}\subseteq\ker\Lambda^{(3)}$.
\end{proposition}

\begin{proof}
At the base multiplicity the constant signed-permutation quotient is trivial
and a pure residual acts as $I+O(h)$.  Since
$\rho_h(W)-I=O(h^3)$ on the whole endpoint module, postcomposition by a pure
residual changes \cref{eq:D-h} only in order $h^4$.  This proves
\cref{eq:beta-cocone} at divided cubic order.  Assumption
\ref{ass:two-handle-cocone} supplies a common physical-copy target and the
same statement at every finite multiplicity.

For a forward pair addition, the two actual core counits act on the new
Frobenius factors.  Under the actual-to-canonical identification in
Assumption~\ref{ass:two-handle-cocone}, the undotted and once-dotted insertions
are $U_b$ and $D_b$.  Lemma~\ref{lem:counit-identities} gives respectively
$0$ and $1$.  The raw degree $+2$ of the once-dotted map cancels the target
cabled shift $-2$, so the second equation has total degree zero.  This proves
\cref{eq:psi-cocone}.  Applying the three equations to the generators in
\cref{eq:two-handle-generators} gives the kernel inclusion.
\end{proof}

\begin{corollary}[Two-handle descent]\label{cor:two-handle-descent}
There is a unique linear map
\begin{equation}\label{eq:Lambda-two-handle}
\overline\Lambda_{2h}^{(3)}:
\mathcal C_{\mathrm{raw}}/\mathcal R_{2h}\longrightarrow\Q
\end{equation}
whose pullback is $\Lambda^{(3)}$.  On the selected class its value is
$2624$.
\end{corollary}

\begin{proof}
Proposition~\ref{prop:beta-psi} places the relation space in the kernel, so
the quotient universal property gives \cref{eq:Lambda-two-handle}.  The
selected base row is the detector of \cref{def:detector}; its value is
\cref{eq:D3-vT}.
\end{proof}

\section{The chosen spheres and three-handle descent}\label{sec:three-handle-descent}

\subsection{The finite homology calculation}

The public data supplies three rooted PL surface constructions.  Their exact
combinatorial sizes are
\begin{center}
\begin{tabular}{@{}lrrr@{}}
\toprule
surface & leaves & pair-of-pants or split bands & root caps\\
\midrule
TH1 & 350176 & 350175 & 1\\
TH2 & 229198 & 229197 & 1\\
THXY & 11115 & 11114 & 1\\
\bottomrule
\end{tabular}
\end{center}
For THXY the band count includes $312$ macro bands and $5401+5401$ corridor
split bands.  The finite checker visits every stored row and verifies the
rooted ancestry, stored normal-framing fields, terminal boundary reduction,
SHA identities, schema, determinant, and verdict strings.  Its geometric
intersection test is narrower: only the comb bands have full coordinates for
an all-row PL self-intersection and slab-gap calculation.  The finger annuli,
bigon, and cap have no such coordinates in the public data.

\begin{computation}[Euler characteristics of the stored surfaces]\label{comp:sphere-euler}
Each of the three stored connected surfaces has Euler characteristic $2$.
\end{computation}

\begin{proof}
The rooted construction begins with one disk for each leaf, attaches one
$1$-handle for each band, and ends with one root cap.  Therefore
\begin{align}
\chi(\mathrm{TH1})&=350176-350175+1=2,\label{eq:chi-th1}\\
\chi(\mathrm{TH2})&=229198-229197+1=2,\label{eq:chi-th2}\\
\chi(\mathrm{THXY})&=11115-11114+1=2.\label{eq:chi-thxy}
\end{align}
The ancestry tree is connected and the orientation data is coherent.  Thus,
within the supplied PL model, each surface is a genus-zero closed oriented
surface.  This calculation does not by itself place the model in the actual
$\partial W_2$; that is Assumption~\ref{ass:chosen-sphere-binding}.
\end{proof}

The stored comb-band coordinate sectors are respectively
\begin{equation}\label{eq:sphere-sectors}
\mathrm{THXY}\subset[2,13/2],\qquad
\mathrm{TH1}\subset[8,9],\qquad
\mathrm{TH2}\subset[10,11].
\end{equation}
The interval separation proves the stated slab-gap property for the stored
comb bands.  It does not prove pairwise disjointness for the complete surfaces.
Embedding the full models as disjoint surfaces in one actual post-two-handle
boundary is Assumption~\ref{ass:chosen-sphere-binding}.

\begin{computation}[Unimodular class matrix]\label{comp:sphere-det}
Let the three claimed sphere classes be the columns of
\begin{equation}\label{eq:sphere-matrix}
S=
\begin{pmatrix}
-1311&-189&41\\
8608&1241&-269\\
-1&0&1
\end{pmatrix}.
\end{equation}
Then $\det S=1$.
\end{computation}

\begin{proof}
Expansion along the third row gives
\begin{align}
\det S
&=(-1)(-40)+1(-39) \label{eq:sphere-det-expand}\\
&=40-39=1. \label{eq:sphere-det-value}
\end{align}
Here the first relevant minor is
$(-189)(-269)-41\cdot1241=-40$, and the third is
$(-1311)1241-(-189)8608=-39$, with the cofactor signs already included in
\cref{eq:sphere-det-expand}.  The same integer calculation is
\path{T73.sphereDet_eq_one} in Lean.
\end{proof}

Under Certified Fact~\ref{comp:finite-spheres} and Assumption~\ref{ass:G4},
the columns of $S$ are the classes of three pairwise-disjoint embedded spheres
in the same $\partial W_2$.  External Result~E7 then permits them to replace the actual attaching
spheres up to the operations in \cref{prop:hj-theorem53}.

\subsection{The six leading maps}

For a chosen sphere $j$, let $b_j>0$ be the number of newly created Frobenius tensor factors.
Use the notation
\begin{equation}\label{eq:UD-def}
D_j=X^{\otimes b_j},\qquad
U_j=\sum_{a=0}^{b_j-1}X^{\otimes a}\otimes1\otimes
X^{\otimes(b_j-1-a)}.
\end{equation}
Let $Q_{s,0}$ and $Q_{t_j,0}$ be the constant terms of the full source and
target coordinate transports.  The abstract leading model is
\begin{align}
F_{j,0}^{(0)}&=Q_{t_j,0}^{-1}
(\id_{\mathrm{old}}\otimes U_j)Q_{s,0},
\label{eq:sphere-undotted-map}\\
F_{j,1}^{(0)}&=Q_{t_j,0}^{-1}
(\id_{\mathrm{old}}\otimes D_j)Q_{s,0}.
\label{eq:sphere-dotted-map}
\end{align}

\begin{lemma}[Abstract sphere rows]\label{lem:abstract-sphere-rows}
Let $r$ be any row on the old factor and put
$E_j=\epsilon^{\otimes b_j}$.  For the maps in
\cref{eq:sphere-undotted-map,eq:sphere-dotted-map},
\begin{equation}\label{eq:abstract-sphere-rows}
(r\otimes E_j)F_{j,0}^{(0)}=0,
\qquad
(r\otimes E_j)F_{j,1}^{(0)}=r,
\end{equation}
after using the displayed coordinate equivalences.
\end{lemma}

\begin{proof}
Coordinate conjugation cancels because the source and target rows are pulled
back through the same $Q_{s,0}$ and $Q_{t_j,0}$.  What remains is
$E_j(U_j)=0$ and $E_j(D_j)=1$, which is
Lemma~\ref{lem:counit-identities}.
\end{proof}

Lemma~\ref{lem:abstract-sphere-rows} is an algebraic theorem.  It does not
identify the abstract maps with the MWW cobordism maps.  That identification,
on the whole source and for the actual chosen surfaces, is exactly
Assumption~\ref{ass:sphere-map-model}.

\begin{proposition}[Three-handle kernel equations]\label{prop:sphere-kernel}
Under Assumption~\ref{ass:G4}, the descended two-handle row $\ell_1$ satisfies,
for each $j\in\{\mathrm{TH1},\mathrm{TH2},\mathrm{THXY}\}$,
\begin{equation}\label{eq:sphere-kernel}
\ell_1\circ\Sigma_{j,0}=0,
\qquad
\ell_1\circ(\Sigma_{j,1}-\id)=0.
\end{equation}
\end{proposition}

\begin{proof}
Assumption~\ref{ass:sphere-map-model} gives the whole-source naturality
square identifying the actual constant terms of $\Sigma_{j,0}$ and
$\Sigma_{j,1}$ with \cref{eq:sphere-undotted-map,eq:sphere-dotted-map}.
The old divided row begins in order $h^3$; the same assumption places every
positive-order transport correction in order at least $h^4$ after
composition.  Hence the divided cubic depends only on the constant maps.
Lemma~\ref{lem:abstract-sphere-rows} then yields the two equations.
\end{proof}

\section{The complete quotient and the conditional conclusion}\label{sec:final-descent}

Let $W_0$ be the one-handle coefficient trace, $W_1$ its quotient by all
two-handle relations, and $W_2$ the quotient of $W_1$ by the six chosen-sphere
relations.  Let $q_{01}:W_0\to W_1$ and $q_{12}:W_1\to W_2$ be the quotient
maps.  Denote the selected class in $W_0$ by $x_0$.

\begin{lemma}[A linear detector survives a quotient]\label{lem:quotient-detector}
Let $V$ be a vector space, $R\subseteq V$ a subspace, and
$\lambda:V\to\Q$ a linear map with $R\subseteq\ker\lambda$.  Then there is a
unique $\bar\lambda:V/R\to\Q$ such that
$\bar\lambda\circ\pi=\lambda$.  If $\lambda(x)\ne0$, then $\pi(x)\ne0$.
\end{lemma}

\begin{proof}
Define $\bar\lambda(x+R)=\lambda(x)$.  If $x-y\in R$, then
$\lambda(x)-\lambda(y)=0$, so the definition is independent of the
representative.  Uniqueness follows from surjectivity of $\pi$.  Finally,
$\pi(x)=0$ would imply $x\in R$ and hence $\lambda(x)=0$.
\end{proof}

\begin{proposition}[Nonzero class after all relation quotients]\label{prop:nonzero-W2}
Under Assumptions~\ref{ass:G2}--\ref{ass:G4},
there is a linear functional $\ell_2:W_2\to\Q$ such that
\begin{equation}\label{eq:ell2-value}
\ell_2(q_{12}q_{01}x_0)=2624.
\end{equation}
In particular, $q_{12}q_{01}x_0\ne0$.
\end{proposition}

\begin{proof}
Lemma~\ref{lem:HH0-descent} first descends the quantum shadow detector through
the coefficient trace.  Proposition~\ref{prop:beta-psi} and
Lemma~\ref{lem:quotient-detector} descend it through $q_{01}$.  Proposition
\ref{prop:sphere-kernel} and the same quotient lemma descend it through
$q_{12}$.  At every stage the pullback agrees with the preceding row, so
\cref{eq:D3-vT} gives \cref{eq:ell2-value}.  Nonvanishing follows from the
last sentence of Lemma~\ref{lem:quotient-detector}.
\end{proof}

\begin{proposition}[Closed candidate class]\label{prop:closed-class}
Under Assumptions~\ref{ass:G4} and \ref{ass:G5}, the class in
\cref{prop:nonzero-W2} maps to a nonzero
class
\begin{equation}\label{eq:closed-class}
z_{494}\in\SLM^2_{0,494}(X(41,189,73);\Q).
\end{equation}
\end{proposition}

\begin{proof}
Assumption~\ref{ass:G4}, through External Result~E4, identifies
$W_2$ with the post-three-handle module after exactly the six displayed
relations.  Assumption~\ref{ass:G5}, through External Result~E5, supplies a
grading-preserving isomorphism after the four-handle attachment.  Isomorphisms preserve
nonzero elements, and Proposition~\ref{prop:degree-494} supplies the quantum
degree.
\end{proof}

\begin{proof}[Proof of Theorem~\ref{thm:main-conditional}]
Proposition~\ref{prop:closed-class} gives a nonzero element in quantum degree
$494$.  By Assumption~\ref{ass:G5}, the same graded piece for $S^4$
is zero.  If the candidate were diffeomorphic to $S^4$, the same assumption,
together with External Result~E6, would give a grading-preserving linear
isomorphism between these two pieces, which is impossible.  Thus the
candidate is not diffeomorphic to $S^4$.  Corollary~\ref{cor:homotopy-sphere}
supplies the homotopy-sphere conclusion.
\end{proof}

\begin{remark}[The $S^4$ control is a zero-to-zero control]\label{rem:s4-control}
When the point-push word is replaced by the identity, the operator
$\rho_h(W)-I$ vanishes and the endpoint detector sends the standard control
to zero.  The published $S^4$ module is also zero in quantum degree $494$.
Thus this control checks compatibility of two zeros; it is not a positive
calibration that independently demonstrates the geometric binding of the
candidate detector.
\end{remark}

\begin{remark}[Presentation and serialization dependence]\label{rem:serialization}
The raw coefficient $[h^3]\ell(\rho_h(W)-I)u$ is a quantity of a based,
labelled, oriented movie presentation.  Reordering the same crossing rows by
the emitter's increasing-coordinate order produces a different series.  With
the corrected collar endpoints its cubic coefficient is $0$ and its first
nonzero coefficient is $663872h^4$.  That order reverses $168$ non-disjoint
events and is not a legal serialization under Assumption
\ref{ass:serialization}.  The former value $-58976$ used the same mixed
endpoint coordinates as the superseded public value.  The calculation
nevertheless shows why
presentation independence cannot be inferred from the scalar alone: it must
be proved through simultaneous transport of every part of the naturality
square.
\end{remark}

\begin{remark}[Insertion-point dependence]\label{rem:insertion-point}
The point-push operator and cap row are inserted at a specified collar and a
specified point.  Moving the insertion point is harmless only if the move is
accompanied by the corresponding trace/cyclicity and collar-isotopy square.
That geometric statement is Assumption~\ref{ass:collar}.  The exact word
calculation does not prove it.
\end{remark}

\section{A fully worked finite-dimensional example}\label{sec:toy-example}

This section illustrates the entire algebraic mechanism on a small example.
It is not a four-manifold calculation and supplies no evidence for any
candidate-specific assumption.  Its purpose is to make the order of
operations---word action, coefficient extraction, relation killing, quotient,
and nonvanishing---checkable by hand.

Let $R_0=\Q[\eps]/(\eps^4)$, $t=1+\eps$, and $E_0=R_0^3$ with basis
$e_0,e_1,e_2$.  The two positive Burau generators are
\begin{align}
B_1&=
\begin{pmatrix}
-\eps&1+\eps&0\\
1&0&0\\
0&0&1
\end{pmatrix},
&
B_2&=
\begin{pmatrix}
1&0&0\\
0&-\eps&1+\eps\\
0&1&0
\end{pmatrix},
\label{eq:toy-positive}\\
B_1^{-1}&=
\begin{pmatrix}
0&1&0\\
t^{-1}&1-t^{-1}&0\\
0&0&1
\end{pmatrix},
&
B_2^{-1}&=
\begin{pmatrix}
1&0&0\\
0&0&1\\
0&t^{-1}&1-t^{-1}
\end{pmatrix},
\label{eq:toy-negative}
\end{align}
where $t^{-1}=1-\eps+\eps^2-\eps^3$.

Take the pure commutator word
\begin{equation}\label{eq:toy-word}
w_0=\sigma_1^2\sigma_2^2\sigma_1^{-2}\sigma_2^{-2}.
\end{equation}
Multiplying the eight matrices from left to right and reducing modulo
$\eps^4$ gives
\begin{equation}\label{eq:toy-product}
\rho(w_0)=
\begin{pmatrix}
1-\eps^3&-\eps^2+\eps^3&\eps^2\\
\eps^2&1&-\eps^2\\
-\eps^2+2\eps^3&\eps^2-3\eps^3&1+\eps^3
\end{pmatrix}.
\end{equation}
The equality can be checked by multiplying
$B_1B_1B_2B_2B_1^{-1}B_1^{-1}B_2^{-1}B_2^{-1}$ using
\cref{eq:toy-positive,eq:toy-negative}; no omitted matrix is involved.

Let $u_0=e_0-e_1$ and let $\ell_0=e_0^*$.  Define the toy detector by the
same pattern as \cref{def:detector}, but at order two:
\begin{equation}\label{eq:toy-detector}
d_2(v)=[\eps^2]\,\ell_0(\rho(w_0)-I)v.
\end{equation}
Reading the coefficient of $\eps^2$ from the first row of
\cref{eq:toy-product} gives the complete detector matrix
\begin{equation}\label{eq:toy-row}
d_2=\begin{pmatrix}0&-1&1\end{pmatrix}.
\end{equation}
Therefore
\begin{equation}\label{eq:toy-selected}
d_2(u_0)=
\begin{pmatrix}0&-1&1\end{pmatrix}
\begin{pmatrix}1\\-1\\0\end{pmatrix}=1.
\end{equation}

Now model two relation maps by the columns
\begin{equation}\label{eq:toy-relations}
r_0=\begin{pmatrix}1\\0\\0\end{pmatrix},
\qquad
r_1=\begin{pmatrix}0\\1\\1\end{pmatrix},
\qquad
R=\begin{pmatrix}1&0\\0&1\\0&1\end{pmatrix}.
\end{equation}
The matrix product
\begin{equation}\label{eq:toy-kernel-check}
d_2R=\begin{pmatrix}0&0\end{pmatrix}
\end{equation}
is the complete kernel check.  Hence $d_2$ descends to
$E_0/\im R$.  A concrete quotient coordinate is
\begin{equation}\label{eq:toy-quotient}
q:E_0\longrightarrow R_0,\qquad
q(x_0,x_1,x_2)=-x_1+x_2,
\end{equation}
whose matrix is the same row as \cref{eq:toy-row}.  The relations vanish:
$q(r_0)=q(r_1)=0$, while $q(u_0)=1$.  Thus the selected class is nonzero in
the quotient.

This toy example contains the exact logical skeleton of the candidate
argument.  What it lacks is precisely the geometric content: there is no
claim that the toy relation columns arise from handle maps or that its final
quotient is a smooth invariant.

\section{Formalization boundary}\label{sec:lean}

The Lean development has two jobs.  First, it checks the finite arithmetic
and the algebra of linear quotients.  Second, it fixes the exact type of the
conditional implication so that a hidden strengthening of the conclusion or
weakening of the premises is difficult.  It does not define smooth
four-manifolds, skein lasagna modules, actual handle movies, or the geometric
predicate that a stored surface is an attaching sphere.

\subsection{What is kernel-checked}

The following statements are proved from definitions in Lean.

\begin{enumerate}
\item The two determinant identities in \cref{comp:determinants}, the chosen
  sphere determinant in \cref{comp:sphere-det}, the integer identity
  $(-2)^3\cdot(-328)=2624$, and the grading sum $-44+227+315-4=494$.
\item The coefficient-$\HH_0$ quotient, the fact that a cyclic shadow kills
  its relation span, and the universal linear factorization through the
  quotient.
\item The Frobenius counit identities
  $\epsilon^{\otimes b}(U_b)=0$ and
  $\epsilon^{\otimes b}(D_b)=1$, including arbitrary positive split trees and
  coordinate conjugation.
\item The fact that six assumed sphere-map bindings place their relation
  ranges in the kernel of the detector, and the induced nonzero-class
  statement in the quotient.
\item The implication from the stated external interfaces to
  ``homotopy sphere and not diffeomorphic to $S^4$.''
\end{enumerate}

The $45{,}360$-letter matrix product is not executed by Lean.  It is executed
by the independent exact Python program described in
\cref{sec:exact-computation}.  Lean imports the final primitive integer
$-328$ for the short substitution check.  This is why the Python computation
and the Lean implication are complementary rather than redundant.

\subsection{What is stored as a field}

The core Lean theorem has the schematic form
\begin{equation}\label{eq:lean-signature}
\texttt{NotStandardInterfaces}\longrightarrow
\neg\texttt{Diffeomorphic(candidate,S4)}.
\end{equation}
The interface contains the actual-$\HH_0$ binding, the actual detector row,
the beta/psi maps, the three actual sphere maps, the MWW coequalizer and
transport, the four-handle map, standard-sphere support, and diffeomorphism
invariance.  Supplying a field of a structure is an assumption, not a proof
of the field's geometric proposition.  In particular, compilation of
\cref{eq:lean-signature} proves only that these fields imply the conclusion.

\small
\begin{longtable}{@{}>{\raggedright\arraybackslash}p{0.24\textwidth}
>{\raggedright\arraybackslash}p{0.18\textwidth}
>{\raggedright\arraybackslash}p{0.49\textwidth}@{}}
\caption{External theorem, paper item, and Lean-field crosswalk.}
\label{tab:crosswalk}\\
\toprule
External result & paper item & principal Lean field or definition\\
\midrule
\endfirsthead
\toprule
External result & paper item & principal Lean field or definition\\
\midrule
\endhead
MWW Theorem 4.7 & E1; used in A2 &
\path{OneHandleInterface.hh0Binding},
\path{hh0Lift}, \path{hh0LiftCommutes}, \path{hh0LiftUnique}\\
BHPW Theorem 4.6 & E9; used in A2 &
candidate-side input to \path{OneHandleInterface.rawCapBinding} and the
naturality fields of \path{SphereChart}\\
BPW Proposition 3.20; Theorem 3.21 & E8; used in A2 &
\path{OneHandleInterface.rawCap},
\path{rawCapKillsHH0}; algebraic quotient in
\path{CoefficientTrace.lean}\\
MWW Definition 3.1; Theorem 3.2 & E2; used in A3 &
\path{BetaPsiInterface.betaRelationBinding},
\path{psi0Binding}, \path{psi1Binding}, \path{q01}\\
Horvat--Jab\l onowski Theorem 5.3 & E7; used in A4 &
\path{SphereMWWFamily.hjBinding} and
\path{SphereTriple.hjBinding}\\
MWW Theorem 3.7; Theorem 3.10 & E3--E4; used in A4 &
\path{SphereMWWFamily.mwwCoequalizerBinding},
\path{q12}, \path{mwwTransport}\\
MWW Proposition 3.4 and Corollary 3.5 & E5; used in A5 &
\path{FourHandleInterface.fourIso},
\path{fourHandleBinding}\\
MWW Definition 5.4 and Example 5.6 & E6; used in A5 &
\path{S4ControlInterface.degreeSupport}; alternatively
\path{S4ReductionData.evalB4} and \path{attach4}\\
Iwaki Proposition 2.1 & E10; used in A1 &
\path{CappellShanesonInterface.matrixConditionsToHomotopySphere}\\
MWW smooth invariance & E6; used in A5 &
\path{DiffeomorphismInvarianceInterface.preservesGradedObject}\\
van Kampen theorem & E11 (alternative route) &
\path{CSTopologyData.vanKampen}\\
Wang and Mayer--Vietoris sequences & E12 (alternative route) &
\path{CSTopologyData.wangMayerVietorisPoincare}\\
Poincar\'e duality & E13 (alternative route) &
\path{CSTopologyData.wangMayerVietorisPoincare}\\
Hurewicz and Whitehead theorems & E14 (alternative route) &
\path{CSTopologyData.hurewiczWhitehead}\\
\bottomrule
\end{longtable}
\normalsize

\subsection{Compilation and axioms}

The recorded clean replay used Lean $4.32.1$ and mathlib commit
\Hash{520045ab14e26149ee970e2e617ca04b09bde5d6}.  It compiled every audited
project module in a fresh output root and then ran
\path{lake lean T73Audit.lean}.  The file requests
\texttt{\#print axioms} for $38$ declarations.  The recorded partition was
two reports using only \texttt{propext}, five using
\texttt{propext, Quot.sound}, and thirty-one using
\texttt{propext, Classical.choice, Quot.sound}.  No report contained
\texttt{sorryAx}.

These are foundational dependencies from mathlib and quotient constructions;
they do not certify the geometric interfaces.  A reader can reproduce both
the compilation and the axiom scan using the commands in
\cref{sec:reproduction}.

\section{Reproduction}\label{sec:reproduction}

The public repository is
\begin{center}
\url{https://github.com/toffee-desuwa/smooth4pc-t73-lean}.
\end{center}
This erratum revision is identified by the repository tag
\texttt{v0.2-erratum1}.
The commands below assume a fresh clone and Python~3.10 or later.

\subsection{Reconstruct the detector}

Run
\begin{lstlisting}
python -I -B scripts/recompute_t73_delta3.py --check
\end{lstlisting}
The expected bearing lines are
\begin{lstlisting}
INPUT_SHA256=0F0F4A568D966367313C29C25CEDDE4FD9ABAD32B9375519E9CDF9168AE0A764
B44_LENGTH=11340
B44_SHA256=7C2D2F792C2672221A76CAF08A71F560AF0CB7654B4D537C00FAD00B16EFA187
B88_LENGTH=45360
B88_SHA256=C13D0A3BA9B05F41A6D2C5B4AB12DDECDABEFC1883835891AD3BD2B8955B5FFB
POSITION_TABLE_SHA256=119C7E9E74AE6C820DA72A84CDFD5D445D81E6C3AACCC209C25D37C323961508
ELL_RHOW_MINUS_I_U_EPS=[0,0,0,-328,14596,-410246,9595271]
ELL_RHOW_MINUS_I_SQUARED_U_EPS=[0,0,0,0,0,0,-102729600]
DELTA3_ETA_T1=2624
DELTA3_XI=0
VERIFY=PASS
\end{lstlisting}
The script itself has SHA-256
\Hash{27B6D1E11CF60FFD13E041B9836A36178EBDA62322FFB046523D55EB67F45CFD}.
The canonical JSON receipt produced from the current input has
SHA-256
\Hash{7C681B99492A5C34B3B8A7EC8B7269E3FD67CD989728C2D4A1C619A63387A722}.

\subsection{Verify the public geometric data}

Run
\begin{lstlisting}
python -I -B scripts/verify_public_geometry_evidence.py
\end{lstlisting}
This command checks the $29$ bundled files, their SHA-256 identities, the
chosen-sphere determinant, and the global-descending calculation from the
full $2{,}126{,}291$-crossing PD input.  Passing this command certifies only
the finite stored model.  It does not prove Assumptions
\ref{ass:actual-presentation}, \ref{ass:chosen-sphere-binding}, or
\ref{ass:sphere-map-model}.

\subsection{Compile and inspect the Lean development}

After installing \texttt{elan}, run
\begin{lstlisting}
lake update
git diff --exit-code -- lake-manifest.json
lake exe cache get
python -B tests/test_t73_minimal_formalization.py -v
lake lean T73Audit.lean
\end{lstlisting}
The Python test requires two passing tests, compiles the audited modules into
a fresh temporary output root, requires exactly $38$ axiom reports, rejects
any axiom outside
\texttt{propext}, \texttt{Classical.choice}, and
\texttt{Quot.sound}, and rejects \texttt{sorryAx}.

\subsection{Build this paper}

On Windows, with Tectonic~0.17.0 installed at the path recorded by
\path{paper/build.ps1}, run
\begin{lstlisting}[language=bash]
powershell -ExecutionPolicy Bypass -File paper/build.ps1
\end{lstlisting}
On Linux or macOS with \texttt{tectonic} on \texttt{PATH}, the cross-platform
one-line command is
\begin{lstlisting}[language=bash]
cd paper && tectonic t73_candidate.tex --outdir build --keep-logs
\end{lstlisting}
The Windows binary used for this draft is Tectonic~0.17.0 with SHA-256
\Hash{99FFCFDBF1EBF8BDDA9E791942E3D06AEDB12463FDDC33F07DE6F5211C8BF08D}.

\section{Review boundary}\label{sec:review-boundary}

The finite computation is an ordinary, reproducible matrix calculation.  Its
nonzero value is not the principal uncertainty.  The principal uncertainty
is whether the stated coefficient, collar, cabled state, and sphere maps are
the actual geometric maps to which the published quotient theorems apply.
The most direct review strategy is therefore not to rerun the last integer
multiplication first.  It is to attack Assumptions~A1--A5 as geometric
statements, in particular the whole-source sphere-map binding in A4.

\begin{remark}[No inference from filenames or certificates]\label{rem:no-certificate-proof}
Names such as ``global descending certificate'', ``chosen sphere receipt'',
or ``Hattori decision'' are locators for data.  They do not occur as premises
of any proof in this article.  Whenever a finite file is used, the relevant
mathematical content is stated explicitly and the replay command is given.
Whenever a geometric conclusion is not derived from those finite data, it is
an Assumption.
\end{remark}

\begin{remark}[The legality audit is bounded]\label{rem:legality-bounded}
The pre-search legality criteria were committed at
\texttt{cf4a990} before the detailed adjudication at
\texttt{9d75dcd}.  Of the $22$ tested alternative standard-member collars,
none was proved legal: $19$ were rejected and $3$ remained unestablished.
This does not prove that every legal standard-member collar has zero detector.
It shows only that the registered t73 chronology survives that bounded test.
\end{remark}

\section{How the five assumptions can be discharged or refuted}\label{sec:assumption-tests}

The five assumptions are not generic requests for ``more justification.''
Each has a definite proof object and a definite failure mode.  This section
states both so that independent review can be adversarial rather than
impressionistic.

\subsection{A1: presentation and collar}

To discharge A1, one must give a framed Kirby movie from the standard
Cappell--Shaneson surgery construction for \cref{eq:A-matrix} to the frozen
cut presentation.  The movie must transport every attaching-component label,
normal framing, basepoint, collar coordinate, and the detector insertion disk.
At each interchange of events, the two supports must be disjoint or an
explicit movie move must be cited.  The final based labelled point-push word
must be the registered word generated in \cref{def:row-expansion,def:cabling}.

A1 is refuted by any forced change of framing, basepoint, insertion point, or
event order that is not accompanied by the simultaneous coordinate transport
of \cref{eq:simultaneous-coordinate}.  It is also refuted if two legal movies
for the same fixed geometric input give different scalar series.  The existing
legality audit rejects the known emitter ordering, but it is not a general
uniqueness theorem for all legal collars.

\subsection{A2: Hattori input and endpoint model}

To discharge A2, one must construct the actual coefficient equivalence
\cref{eq:H-balanced} from the paired annuli and prove both action-naturality
squares.  Applying the quantum trace and the $227$ counits must then give the
specific vector $u=e_2-e_{87}$, with the covector
$\ell=e_{87}^*-e_2^*$ arising from the same collar position table.  The proof
must also bind the actual word to the third pure-braid filtration term on the
whole endpoint module and derive the four grading contributions from that
same filling.

A2 is refuted if the actual diagonal class maps to a different endpoint
vector, if the cap row changes without simultaneous conjugation, if the word
has a lower-order term on any endpoint vector, or if one grading contribution
belongs to a different normalization.  Agreement of the single scalar
$2624$ is not enough to rule out these failures.

\subsection{A3: the all-state two-handle cocone}

To discharge A3, one must define the statewise rows for every finite cable
state and prove \cref{eq:beta-cocone,eq:psi-cocone} as equalities of linear
maps on their full source spaces.  The common-target and sign choices must be
compatible around every square of the cable-state graph.  A proof on the
selected vector, a generic-rank computation, or a finite list of cable states
does not establish A3.

A3 is refuted by a single admissible braid or stabilization source vector on
which one of those map equalities fails.  Such a counterexample would show
that the detector does not descend through the MWW two-handle quotient even
though its value on the base representative is nonzero.

\subsection{A4: actual spheres and whole-source sphere maps}

To discharge A4, the finite surfaces of C2 must first be embedded in one
actual $\partial W_2$ with their framings, disjointness, and homology classes
preserved.  The hypotheses of External Result E7 must then be verified for
that same handle decomposition.  Finally, the two MWW cobordism maps for each
sphere must be identified, on the whole source, with the conjugated leading
maps \cref{eq:sphere-undotted-map,eq:sphere-dotted-map}; their positive-order
transport terms must be shown not to enter the divided cubic.

A4 is refuted if any chosen surface fails to bind to the actual boundary, if
the replacement theorem's boundary decomposition or basis hypothesis fails,
or if one actual sphere-map vector has a nonzero detector value contrary to
\cref{eq:sphere-kernel}.  Because this is the last relation layer before the
closed invariant, A4 is the most consequential unresolved point.

\subsection{A5: rational coefficients and grading}

To discharge A5, one must trace the coefficient functor from the integral
$B^4$ evaluation and $S^4$ computation to the rational module used here,
showing that tensoring introduces neither a Tor correction nor a grading
change.  The four-handle isomorphism must be checked to have bidegree $(0,0)$,
and the diffeomorphism-induced equivalence must preserve the same quantum
grading.

A5 is refuted if the rational degree-$494$ piece of $S^4$ is not zero in the
chosen convention, or if either comparison map shifts degree.  Unlike A3 and
A4, this question is largely a matter of matching published coefficient and
grading conventions, so it should be settled before undertaking further large
finite computations.

\appendix

\section{The complete degree-three endpoint vector}\label{app:vector}

With the conventions of \cref{sec:endpoint}, the coefficient vector in
\cref{eq:eps3-vector-short} is
\begin{align}
(a_0,\ldots,a_{87})={}&(
-7052,-7052,164,164,332,332,324,324,316,316,308,308,\notag\\
&300,300,292,292,284,284,276,276,268,268,260,260,\notag\\
&252,252,244,244,236,236,228,228,220,220,212,212,\notag\\
&204,204,196,196,188,188,180,180,172,172,164,164,\notag\\
&156,156,148,148,140,140,132,132,124,124,116,116,\notag\\
&108,108,100,100,92,92,84,84,76,76,68,68,\notag\\
&60,60,52,52,44,44,36,36,28,28,20,20,\notag\\
&12,12,-164,-164).
\label{eq:full-88-vector}
\end{align}
The detector covector is $\ell=e_{87}^*-e_2^*$, so the contraction is visibly
\begin{equation}\label{eq:vector-contraction}
\ell(a_0,\ldots,a_{87})^T=a_{87}-a_2=-164-164=-328.
\end{equation}
The sum of the $88$ entries is zero, as expected for the unreduced Burau
difference on the invariant total-coordinate functional; this is a check, not
an input to \cref{eq:vector-contraction}.

\section{The four grading contributions}\label{app:grading}

For ease of audit, the four contributions of \cref{sec:grading} are collected
without abbreviation:
\begin{longtable}{@{}p{0.17\textwidth}p{0.45\textwidth}p{0.27\textwidth}@{}}
\toprule
term & source & calculation\\
\midrule
$-44$ & remove the $p(N-1)$ Hom normalization for $p=44$, $N=2$ & $-44(2-1)=-44$\\
$+227$ & $227$ distinct Frobenius labels, each of degree $+1$ & $227\cdot1=227$\\
$+315$ & one-handle gluing shift from $44$ $y$-pairs and $271=44+227$ $z$-pairs & $44+271=315$\\
$-4$ & cabled shift for $N=2$, $\alpha=0$, $|r|=2$ & $(1-2)(2\cdot2+0)=-4$\\
\midrule
total & & $-44+227+315-4=494$\\
\bottomrule
\end{longtable}

\section{Convention-legality criteria}\label{app:legality}

A serialization is legal only if it describes the same framed handle
presentation, the same oriented based cut tangle, the same labelled boundary
pairing, the same normal field, and the same insertion point.  Its event order
must be a linear extension of the geometric partial order.  Adjacent events
may be exchanged only when their supports are disjoint and a strict
interchange theorem applies.  Word replacement requires equality in the
based labelled braid or tangle category; equality of closures, permutations,
free-group words, or homology classes is insufficient.

A coordinate change $P$ must act simultaneously:
\begin{equation}\label{eq:simultaneous-coordinate}
W'=PWP^{-1},\qquad u'=Pu,\qquad
\ell'=\ell P^{-1},
\end{equation}
with the shadow, collar, basepoint, and insertion point transported in the
same naturality square.  Under \cref{eq:simultaneous-coordinate},
\begin{equation}\label{eq:coordinate-invariance}
\ell'(W'-I)u'=\ell P^{-1}P(W-I)P^{-1}Pu=\ell(W-I)u.
\end{equation}
Thus coordinate invariance is proved when the simultaneous transport exists;
it is not imposed by selecting the preferred numerical answer.

A collar is legal only if it is an actual framed tubular collar of the fixed
component and is isotopic to the registered collar relative to the boundary,
basepoints, labels, old input balls, normal framing, and insertion point.  A
nontrivial relative mapping-class element or an added framing twist changes
the input.  These criteria were written before the detailed alternative
collar outputs were opened.  Their use in the candidate geometry remains
Assumption~\ref{ass:G1}.

\section{Terminology}\label{app:glossary}

\begin{description}
\item[Actual map]
A map induced by a specified geometric tangle or surface cobordism, as
opposed to an abstract linear map having the desired source and target.
\item[Candidate-specific binding]
An identification between immutable finite data and the actual geometric
object required by a published theorem.
\item[Complete quotient]
The quotient by all relation generators in the applicable MWW handle formula,
not by a sampled finite subset.
\item[Divided cubic]
The reduction modulo $h$ of a map known on its whole domain to be divisible by
$h^3$, as in \cref{eq:D3}.
\item[Endpoint shadow]
The image of a coefficient trace class under the quantum
vertical-to-horizontal trace and the fixed endpoint functor.
\item[Finite certificate]
Immutable finite data plus a deterministic exact checker.  The term has no
geometric force beyond the predicates actually recomputed.
\item[Registered movie order]
The oriented point-push chronology fixed in the public primitive input.  In
standard mathematical language it is the chosen linear extension of the
geometric event partial order.
\item[Material factor]
In the earlier project notes this meant a newly created Frobenius tensor
factor.  The paper uses the latter standard phrase.
\item[Reynolds average]
The finite group average over physical-copy permutations.  It is used only in
the assumed all-state beta/psi construction, not as an independent geometric
theorem.
\end{description}

\bibliographystyle{amsplain}
\bibliography{references}

\end{document}
